\documentclass[a4paper,12pt,openany,oneside]{book}
\usepackage[utf8]{inputenc}
\usepackage{multicol}
\usepackage{amsmath}
\usepackage{amssymb}
\usepackage{amsthm}
\usepackage{geometry}
\usepackage{graphicx}
\usepackage[hidelinks]{hyperref}
\usepackage{listings}
\usepackage{xcolor}
\usepackage{algorithm}
\usepackage{algpseudocode}
\usepackage{chngcntr}
\usepackage[protrusion=true,expansion=true,final,babel=false]{microtype}
\usepackage{enumitem}
\usepackage{booktabs}
\usepackage{array}
\usepackage{ragged2e}
\usepackage{caption}

\geometry{a4paper,left=2.8cm,right=2.8cm,top=2.5cm,bottom=2.8cm,
          headheight=0pt,headsep=0pt}
\pagestyle{plain}

\hypersetup{colorlinks=true,linkcolor=blue!60!black,
            urlcolor=blue!60!black,citecolor=blue!60!black}

%--- Algorithmus-Bezeichner ---
\floatname{algorithm}{Algorithmus}
\renewcommand{\listalgorithmname}{Algorithmenverzeichnis}
\renewcommand{\algorithmicrequire}{\textbf{Eingabe:}}
\renewcommand{\algorithmicensure}{\textbf{Ausgabe:}}
\renewcommand{\algorithmicwhile}{\textbf{solange}}
\renewcommand{\algorithmicdo}{\textbf{tue}}
\renewcommand{\algorithmicend}{\textbf{Ende}}
\renewcommand{\algorithmicif}{\textbf{falls}}
\renewcommand{\algorithmicthen}{\textbf{dann}}
\renewcommand{\algorithmicelse}{\textbf{sonst}}
\renewcommand{\algorithmicfor}{\textbf{fuer}}
\renewcommand{\algorithmicforall}{\textbf{fuer alle}}
\renewcommand{\algorithmicreturn}{\textbf{Rueckgabe}}
\renewcommand{\algorithmicrepeat}{\textbf{wiederhole}}
\renewcommand{\algorithmicuntil}{\textbf{bis}}
\renewcommand{\algorithmiccomment}[1]{\hfill\textcolor{gray}{$\triangleright$ #1}}

%--- Theorem-Umgebungen ---
\theoremstyle{definition}
\newtheorem{definition}{Definition}[chapter]
\newtheorem{beispiel}{Beispiel}[chapter]
\newtheorem{frage}{Frage}[chapter]
\newtheorem{aufgabe}{Aufgabe}[chapter]
\theoremstyle{plain}
\newtheorem{satz}{Satz}[chapter]
\newtheorem{lemma}{Lemma}[chapter]
\newtheorem{korollar}{Korollar}[chapter]
\newtheorem{proposition}{Proposition}[chapter]
\theoremstyle{remark}
\newtheorem{bemerkung}{Bemerkung}[chapter]

\counterwithin{algorithm}{chapter}
\counterwithin{figure}{chapter}
\counterwithin{table}{chapter}

%--- Listings ---
\lstset{basicstyle=\ttfamily\footnotesize,
  keywordstyle=\color{blue!80!black}\bfseries,
  commentstyle=\color{gray!70!black}\itshape,
  stringstyle=\color{red!70!black},
  numbers=left,numberstyle=\tiny\color{gray},
  breaklines=true,breakatwhitespace=true,
  frame=single,tabsize=4,showstringspaces=false,
  captionpos=b,xleftmargin=1.5em,xrightmargin=0.5em}
\renewcommand{\lstlistingname}{Quelltext}
\renewcommand{\lstlistlistingname}{Quelltextverzeichnis}

%--- Listen ---
\setlist[itemize,1]{label=$-$,leftmargin=1.5em,topsep=3pt,itemsep=1pt}
\setlist[enumerate,1]{leftmargin=1.5em,topsep=3pt,itemsep=1pt}

\sloppy
\setlength{\emergencystretch}{3em}


\title{\Huge\textbf{Graphentheorie}\\[1.2ex]
  \Large Eine algorithmisch-orientierte Einf\"uhrung\\[0.8ex]
  \large mit Graphprofilen, Signatur-Methode und Anwendungen}
\author{Stephan Epp\\[0.5ex]\small Universit\"at Bielefeld}
\date{\today}






% Backmatter-Überschrift ohne erzwungenen Seitenumbruch
\newcommand{\backchapter}[1]{%
  \phantomsection%
  \addcontentsline{toc}{chapter}{#1}%
  {\normalfont\huge\bfseries #1\par}%
  \vspace*{20pt}%
  \nobreak%
}

% Unterdrückt den internen \chapter*-Aufruf von thebibliography
\makeatletter
\renewenvironment{thebibliography}[1]{%
  \list{\@biblabel{\@arabic\c@enumiv}}{%
    \settowidth\labelwidth{\@biblabel{#1}}%
    \leftmargin\labelwidth
    \advance\leftmargin\labelsep
    \@openbib@code
    \usecounter{enumiv}%
    \let\p@enumiv\@empty
    \renewcommand\theenumiv{\@arabic\c@enumiv}}%
  \sloppy
  \clubpenalty4000
  \@clubpenalty\clubpenalty
  \widowpenalty4000%
  \sfcode`\.\@m}
  {\def\@noitemerr
    {\@latex@warning{Empty `thebibliography' environment}}%
   \endlist}
\makeatother
\AtBeginDocument{%
  \renewcommand{\contentsname}{Inhaltsverzeichnis}%
  \renewcommand{\bibname}{Literaturverzeichnis}%
  \renewcommand{\chaptername}{Kapitel}%
  \renewcommand{\figurename}{Abbildung}%
  \renewcommand{\tablename}{Tabelle}%
  \renewcommand{\listfigurename}{Abbildungsverzeichnis}%
  \renewcommand{\listtablename}{Tabellenverzeichnis}%
  \renewcommand{\indexname}{Stichwortverzeichnis}%
  \renewcommand{\appendixname}{Anhang}%
  \renewcommand{\partname}{Teil}%
}
\begin{document}
% Deutsche Strukturbegriffe
\renewcommand{\proofname}{Beweis}
\renewcommand{\chaptername}{Kapitel}
\renewcommand{\appendixname}{Anhang}
\renewcommand{\contentsname}{Inhaltsverzeichnis}
\renewcommand{\listfigurename}{Abbildungsverzeichnis}
\renewcommand{\listtablename}{Tabellenverzeichnis}
\renewcommand{\bibname}{Literaturverzeichnis}
\renewcommand{\indexname}{Stichwortverzeichnis}
\renewcommand{\figurename}{Abbildung}
\renewcommand{\tablename}{Tabelle}
\renewcommand{\partname}{Teil}
\renewcommand{\today}{%
  \number\day.~\ifcase\month\or
  Januar\or Februar\or M\"arz\or April\or Mai\or Juni\or
  Juli\or August\or September\or Oktober\or November\or Dezember%
  \fi\space\number\year}
% KOMA-spezifische Bezeichner
\frontmatter
\maketitle
\clearpage

\renewcommand{\contentsname}{Inhaltsverzeichnis}
\tableofcontents
\clearpage
\chapter*{Vorwort}
\addcontentsline{toc}{chapter}{Vorwort}

Graphen geh\"oren zu den m\"achtigsten und vielseitigsten Abstraktionen
der Informatik und Mathematik. Ein Graph besteht aus einer Menge von
Knoten und einer Menge von Kanten, die Paare von Knoten verbinden.
Diese schlichte Struktur gen\"ugt, um eine au\ss{}erordentlich breite
Klasse realer Systeme zu modellieren: Stra\ss{}ennetze, Computernetzwerke,
soziale Beziehungsgeflechte, Proteininteraktionen, logische Abh\"angigkeiten
in Software sowie -- besonders eindrucksvoll -- die synaptischen Verbindungen
zwischen Neuronen im menschlichen Gehirn.

Die St\"arke des Graphmodells liegt in seiner Abstraktion. Sobald ein
System als Graph formalisiert ist, lassen sich allgemeine algorithmische
Werkzeuge anwenden, unabh\"angig vom jeweiligen Anwendungsgebiet. Fragen
wie \emph{Wie weit sind zwei Knoten voneinander entfernt?}, \emph{Gibt es
einen Weg von $A$ nach $B$?} oder \emph{Welche Knoten sind f\"ur die
Konnektivit\"at besonders kritisch?} \"ubersetzen sich direkt in
algorithmische Probleme auf Graphen.

Das Besondere dieser Einf\"uhrung ist die konsequente \emph{algorithmische
Perspektive}. Neben der klassischen Graphentheorie -- entwickelt entlang
der Kapitelstruktur eines anerkannten Lehrbuches -- wird die
\emph{Graphprofilverteilung} als zentrales Analysewerkzeug eingef\"uhrt.
Das \emph{Profil} eines Graphen -- das Tripel aus k\"urzester-Wege-Matrix
$D$, l\"angster-Wege-Matrix $L$ und Kantenma\ss{} $\kappa$ -- erlaubt eine
optimale, deterministische Charakterisierung jedes Graphen und bildet die
Grundlage f\"ur den Vergleich von Netzwerken verschiedenster Dom\"anen.

Ein weiteres Alleinstellungsmerkmal ist die \emph{Signatur-Methode} f\"ur
die Boolean Matrixmultiplikation. Diese Technik, bei der jede Zeile einer
Boolean-Matrix als ganzzahlige Signatur kodiert wird, erm\"oglicht eine
Laufzeit von $O(n^2)$ f\"ur die Matrixmultiplikation und damit $O(n^3)$
f\"ur die vollst\"andige Profilberechnung -- jeweils bei $O(n^2)$ Speicher.
Alle Algorithmen werden mit vollst\"andigen Korrektheits- und
Laufzeitnachweisen pr\"asentiert.

\medskip
\noindent\textbf{Aufbau des Buches.}
Kapitel~1 legt die Grundbegriffe (Graphen, Wege, Zusammenhang) fest.
Kapitel~2 f\"uhrt die algebraischen Darstellungen (Adjazenzmatrix,
Laplace-Matrix, Kantenma\ss{}) ein und beweist den Zusammenhang zwischen
Matrixpotenzen und Wegen.
Kapitel~3 behandelt planare Graphen und die Eulersche Polyederformel.
Kapitel~4 untersucht unabh\"angige Mengen, Cliquen und Matchings.
Kapitel~5 widmet sich der Graphf\"arbung inklusive des Vier-Farben-Satzes
und des chromatischen Polynoms.
Kapitel~6 analysiert den Zusammenhang von Graphen quantitativ.
Kapitel~7 behandelt B\"aume, die Cayleysche Formel und den Pr\"ufercode.
Kapitel~8 untersucht Kreise: Eulerkreise, Hamiltonkreise und den Zyklenraum.
Kapitel~9 f\"uhrt gerichtete Graphen mit Anwendungen in Fl\"ussen und
neuronalen Netzwerken ein.
Kapitel~10 behandelt spezielle Gebiete: extremale Graphentheorie,
Symmetrie, chordale Graphen und dominierende Mengen.
Kapitel~11 -- das Kernkapitel -- pr\"asentiert die Graphenalgorithmen
(DFS, BFS, Dijkstra, Floyd-Warshall, vollst\"andige Profilberechnung) mit
s\"amtlichen Korrektheits- und Laufzeitnachweisen sowie die
Graphprofilverteilung mit experimentellen Daten.
Kapitel~12 gibt einen ausf\"uhrlichen Ausblick auf offene Forschungsfragen.

\medskip
\noindent\textbf{Voraussetzungen.}
Das Buch setzt mathematische Grundkenntnisse voraus, wie sie in den
ersten Studiensemestern der Informatik oder Mathematik vermittelt werden:
elementare Mengenlehre, vollst\"andige Induktion, Grundlagen der
Linearen Algebra. Programmierkenntnisse in Python sind f\"ur die
Quelltextbeispiele hilfreich, aber nicht zwingend erforderlich.

\bigskip
\noindent Bielefeld, \today

\hfill Stephan Epp

\mainmatter
\chapter{Graphen -- Grundbegriffe und Strukturen}

Graphen geh\"oren zu den m\"achtigsten und vielseitigsten Abstraktionen der
Informatik und Mathematik. Ein Graph besteht aus einer Menge von Knoten und
einer Menge von Kanten, die Paare von Knoten verbinden. Diese schlichte
Struktur gen\"ugt, um eine au\ss{}erordentlich breite Klasse realer Systeme
zu modellieren: Stra\ss{}ennetze, Computernetzwerke, soziale Beziehungsgeflechte,
Proteininteraktionen, logische Abh\"angigkeiten in Software sowie -- besonders
eindrucksvoll -- die synaptischen Verbindungen zwischen Neuronen im menschlichen
Gehirn.

Die St\"arke des Graphmodells liegt in seiner Abstraktion. Sobald ein System
als Graph formalisiert ist, lassen sich allgemeine algorithmische Werkzeuge
anwenden, unabh\"angig vom jeweiligen Anwendungsgebiet. Fragen wie \emph{Wie
weit sind zwei Knoten voneinander entfernt?}, \emph{Gibt es einen Weg von $A$
nach $B$?} oder \emph{Welche Knoten sind f\"ur die Konnektivit\"at besonders
kritisch?} \"ubersetzen sich direkt in algorithmische Probleme auf Graphen,
deren L\"osung in allen Anwendungsdom\"anen gleicherma\ss{}en n\"utzlich ist.

\section{Grundlegende Definitionen}

\begin{definition}[Graph]
  Ein \emph{(ungerichteter) Graph} $G = (V, E)$ besteht aus einer endlichen
  nichtleeren \emph{Knotenmenge} $V$ und einer \emph{Kantenmenge}
  $E \subseteq \bigl\{\{u,v\} \mid u,v \in V,\; u \neq v\bigr\}$.
  Die Anzahl der Knoten $n = |V|$ hei\ss{}t \emph{Ordnung}, die Anzahl der
  Kanten $m = |E|$ hei\ss{}t \emph{Gr\"o\ss{}e} des Graphen. Wir schreiben auch
  $V(G)$ f\"ur die Knotenmenge und $E(G)$ f\"ur die Kantenmenge von $G$.
\end{definition}

\begin{bemerkung}
  Die Bezeichnungen $V$ und $E$ entstammen den englischen W\"ortern
  \emph{vertex} (Knoten) und \emph{edge} (Kante). In der gesamten Arbeit
  verwenden wir $n = |V|$ und $m = |E|$ als Standardnotation. Wenn wir von
  einem Graphen ohne weitere Zusatz sprechen, meinen wir stets einen
  endlichen, schlichten, ungerichteten Graphen.
\end{bemerkung}

\begin{definition}[Schlinge, Parallele Kanten, Schlichter Graph, Multigraph]
  Eine Kante $\{v,v\}$, deren beide Endknoten zusammenfallen, hei\ss{}t
  \emph{Schlinge}. Zwei Kanten, die dieselben Endknoten verbinden, hei\ss{}en
  \emph{parallele Kanten}. Ein Graph ohne Schlingen und ohne parallele Kanten
  hei\ss{}t \emph{schlichter Graph}. Ein Graph, der parallele Kanten oder
  Schlingen enthalten kann, hei\ss{}t \emph{Multigraph}. In dieser Arbeit
  betrachten wir -- sofern nicht anders angegeben -- ausschlie\ss{}lich
  schlichte Graphen.
\end{definition}

\begin{definition}[Gerichteter Graph, Digraph]
  Ein \emph{gerichteter Graph} (kurz \emph{Digraph}) $G = (V, E)$ besitzt
  eine Kantenmenge $E \subseteq V \times V$, bei der jede Kante $(u,v)$ eine
  feste Richtung von $u$ (dem \emph{Startknoten}) nach $v$ (dem
  \emph{Zielknoten}) tr\"agt. Eine Kante $(u,v)$ in einem Digraphen hei\ss{}t
  auch \emph{Pfeil} oder \emph{gerichtete Kante}.
\end{definition}

\subsection{Knotengrade}

\begin{definition}[Grad, Minimalgrad, Maximalgrad]
  Der \emph{Grad} $\deg_G(v)$ eines Knotens $v \in V$ ist die Anzahl der
  Kanten in $E$, die $v$ als Endknoten besitzen:
  \[
    \deg_G(v) = \bigl|\{ e \in E \mid v \in e \}\bigr|.
  \]
  Wir schreiben kurz $\deg(v)$, wenn der Graph $G$ aus dem Kontext klar ist.
  Der \emph{Minimalgrad} ist $\delta(G) = \min_{v \in V} \deg(v)$ und der
  \emph{Maximalgrad} ist $\Delta(G) = \max_{v \in V} \deg(v)$.
  Ein Knoten mit $\deg(v) = 0$ hei\ss{}t \emph{isolierter Knoten}; ein Knoten
  mit $\deg(v) = 1$ hei\ss{}t \emph{Blatt}.
\end{definition}

\begin{definition}[Nachbarschaft]
  Die \emph{offene Nachbarschaft} eines Knotens $v$ ist die Menge seiner
  Nachbarknoten: $N(v) = \{ u \in V \mid \{u,v\} \in E \}$. Die
  \emph{geschlossene Nachbarschaft} ist $N[v] = N(v) \cup \{v\}$.
  Es gilt $\deg(v) = |N(v)|$.
\end{definition}

\begin{satz}[Handshaking-Lemma]
  \label{satz:handshaking}
  F\"ur jeden Graphen $G = (V, E)$ gilt:
  \[
    \sum_{v \in V} \deg(v) = 2\,|E| = 2m.
  \]
  Insbesondere ist die Anzahl der Knoten mit ungeradem Grad stets gerade.
\end{satz}

\begin{proof}
  Jede Kante $\{u,v\} \in E$ hat genau zwei Endknoten. Daher tr\"agt sie zur
  Summe $\sum_{v} \deg(v)$ genau zweimal bei -- einmal f\"ur $u$ und einmal
  f\"ur $v$. Somit gilt $\sum_{v} \deg(v) = 2m$. Da $2m$ eine gerade Zahl ist
  und die Summe aller Grade gerade ist, muss die Anzahl der Summanden, die
  ungerade sind (d.\,h. die Anzahl der Knoten mit ungeradem Grad), ebenfalls
  gerade sein.
\end{proof}

\begin{korollar}
  In jedem Graphen ist die Anzahl der Knoten mit ungeradem Grad gerade.
  Insbesondere hat jeder $k$-regul\"are Graph (alle Grade gleich $k$) die
  Eigenschaft $k \cdot n = 2m$; eine solche Aussage l\"asst sich nicht mit
  $k$ ungerade und $n$ ungerade erf\"ullen.
\end{korollar}

\begin{beispiel}[Gradfolgen]
  Wir betrachten einen Graphen $G$ mit Knotenmenge $V = \{1,2,3,4,5\}$ und
  Kantenmenge $E = \{\{1,2\},\{1,3\},\{2,3\},\{3,4\},\{4,5\}\}$. Dann ist
  $m = 5$ und die Grade sind: $\deg(1) = 2$, $\deg(2) = 2$, $\deg(3) = 3$,
  $\deg(4) = 2$, $\deg(5) = 1$. Die Summe $2+2+3+2+1 = 10 = 2 \cdot 5$
  best\"atigt das Handshaking-Lemma.
\end{beispiel}

\begin{definition}[Gradfolge, gradfolgenartig]
  Die \emph{Gradfolge} eines Graphen ist die monoton nicht-wachsend sortierte
  Folge aller Knotengrade: $d_1 \geq d_2 \geq \cdots \geq d_n$.
  Eine Folge nichtnegativer ganzer Zahlen $(d_1, \ldots, d_n)$ mit
  $d_1 \geq \cdots \geq d_n$ hei\ss{}t \emph{graphisch}, wenn es einen
  Graphen mit genau dieser Gradfolge gibt.
\end{definition}

\begin{satz}[Erd\H{o}s--Gallai, 1960]
  Eine Folge $(d_1, d_2, \ldots, d_n)$ mit $d_1 \geq d_2 \geq \cdots \geq
  d_n \geq 0$ ist genau dann graphisch, wenn $\sum_{i=1}^n d_i$ gerade ist
  und f\"ur alle $k \in \{1,\ldots,n\}$ gilt:
  \[
    \sum_{i=1}^k d_i \leq k(k-1) + \sum_{i=k+1}^n \min(d_i, k).
  \]
\end{satz}

\subsection{Wege und Kreise}

\begin{definition}[Weg, einfacher Weg, L\"ange]
  Ein \emph{Weg} $W$ der L\"ange $k$ von Knoten $u$ nach Knoten $v$ ist eine
  endliche Folge $(v_0, v_1, \ldots, v_k)$ von Knoten mit $v_0 = u$,
  $v_k = v$ und $\{v_l, v_{l+1}\} \in E$ f\"ur alle $0 \leq l < k$.
  Der Weg hei\ss{}t \emph{einfach}, wenn alle Knoten $v_0, v_1, \ldots, v_k$
  paarweise verschieden sind. Die \emph{L\"ange} eines Weges ist die Anzahl
  seiner Kanten.
\end{definition}

\begin{definition}[Kreis, einfacher Kreis, Taille, Umfang]
  Ein \emph{Kreis} ist ein geschlossener Weg $(v_0, v_1, \ldots, v_k)$ mit
  $v_0 = v_k$ und $k \geq 1$. Ein Kreis hei\ss{}t \emph{einfach}, wenn alle
  Knoten $v_0, v_1, \ldots, v_{k-1}$ paarweise verschieden sind. Die
  \emph{Taille} $\mathrm{g}(G)$ eines Graphen ist die L\"ange eines k\"urzesten
  Kreises (f\"ur kreisfreie Graphen: $\mathrm{g}(G) = \infty$), der
  \emph{Umfang} $\mathrm{circ}(G)$ ist die L\"ange eines l\"angsten Kreises.
  Ein Graph ohne Kreise hei\ss{}t \emph{azyklisch} oder \emph{Wald}.
\end{definition}

\begin{bemerkung}
  Der Begriff \emph{Weg} wird in der Literatur uneinheitlich verwendet.
  Manche Autoren (darunter Tittmann~\cite{tittmann2025}) unterscheiden
  zwischen \emph{Weg} (alle Kanten verschieden) und \emph{einfachem Weg}
  (alle Knoten und Kanten verschieden). Wir verwenden im Folgenden
  \emph{Weg} stets im Sinne eines einfachen Weges, sofern nicht anders
  angegeben.
\end{bemerkung}

\subsection{Zusammenhang}

\begin{definition}[Zusammenhang, Zusammenhangskomponenten]
  Zwei Knoten $u, v \in V$ hei\ss{}en \emph{verbunden}, wenn es einen Weg von
  $u$ nach $v$ in $G$ gibt. Ein Graph $G$ hei\ss{}t \emph{zusammenh\"angend},
  wenn je zwei seiner Knoten verbunden sind. Andernfalls hei\ss{}t $G$
  \emph{unzusammenh\"angend}. Die maximalen zusammenh\"angenden Teilgraphen
  von $G$ hei\ss{}en \emph{Zusammenhangskomponenten} von $G$.
\end{definition}

\begin{definition}[Br\"ucke, Artikulation]
  Eine Kante $e \in E$ hei\ss{}t \emph{Br\"ucke}, wenn $G - e$ (der Graph mit
  entfernter Kante $e$) mehr Zusammenhangskomponenten als $G$ hat. Ein
  Knoten $v \in V$ hei\ss{}t \emph{Artikulation}, wenn $G - v$ (der Graph mit
  entferntem Knoten $v$ und allen inzidenten Kanten) mehr
  Zusammenhangskomponenten als $G$ hat.
\end{definition}

\begin{satz}
  \label{satz:brucke-kreis}
  Eine Kante $e$ ist genau dann eine Br\"ucke, wenn sie in keinem Kreis von
  $G$ enthalten ist.
\end{satz}

\begin{proof}
  ($\Rightarrow$) Sei $e = \{u,v\}$ eine Br\"ucke. Dann gilt: Wenn es einen
  Kreis $C$ durch $e$ g\"abe, so g\"abe es auch einen Weg von $u$ nach $v$
  in $G - e$ (n\"amlich den Weg entlang $C$ ohne die Kante $e$). Dies
  widerspricht der Tatsache, dass $e$ eine Br\"ucke ist.
  ($\Leftarrow$) Sei $e = \{u,v\}$ in keinem Kreis enthalten. Dann ist der
  einzige Weg von $u$ nach $v$ in $G$ die Kante $e$ selbst. Nach Entfernen
  von $e$ sind $u$ und $v$ nicht mehr verbunden, also hat $G - e$ mehr
  Zusammenhangskomponenten als $G$.
\end{proof}

\section{Operationen mit Graphen}

\begin{definition}[Teilgraph, Induzierter Teilgraph, Spannender Teilgraph]
  Ein Graph $H = (V', E')$ hei\ss{}t \emph{Teilgraph} (Subgraph) von
  $G = (V,E)$, falls $V' \subseteq V$ und $E' \subseteq E$. $H$ hei\ss{}t
  \emph{induzierter Teilgraph} von $G$ durch $V'$, geschrieben $G[V']$, falls
  $E' = \{\{u,v\} \in E \mid u,v \in V'\}$, d.\,h. alle Kanten aus $G$
  zwischen Knoten von $V'$ werden beibehalten. $H$ hei\ss{}t
  \emph{spannender Teilgraph}, wenn $V' = V$.
\end{definition}

\begin{definition}[Komplement]
  Das \emph{Komplement} $\overline{G}$ eines schlichten Graphen $G = (V,E)$
  hat dieselbe Knotenmenge $V$ und als Kantenmenge alle Paare, die in $G$
  fehlen:
  $E(\overline{G}) = \bigl\{\{u,v\} \mid u \neq v, \{u,v\} \notin E\bigr\}.$
  Es gilt $G \cup \overline{G} = K_n$ (vollst\"andiger Graph).
\end{definition}

\begin{definition}[Kontraktion]
  Die \emph{Kontraktion} $G / e$ einer Kante $e = \{u,v\}$ entsteht, indem
  $u$ und $v$ zu einem neuen Knoten $w$ verschmolzen werden und $w$ mit allen
  fr\"uheren Nachbarn von $u$ und $v$ verbunden wird (au\ss{}er $u$ und $v$
  selbst).
\end{definition}

\section{Spezielle Graphklassen}

\begin{definition}[Vollst\"andiger Graph $K_n$]
  Der \emph{vollst\"andige Graph} $K_n$ hat $n$ Knoten und enth\"alt alle
  $\binom{n}{2} = \frac{n(n-1)}{2}$ m\"oglichen Kanten. Jeder Knoten hat
  Grad $n-1$. Das Kantenma\ss{} (vgl.~Kapitel~\ref{chap:matrizen}) betr\"agt
  $\kappa(K_n) = \frac{2}{n-1} \to 0$ f\"ur $n \to \infty$.
\end{definition}

\begin{definition}[Pfadgraph $P_n$, Kreisgraph $C_n$]
  Der \emph{Pfadgraph} $P_n$ hat $n$ Knoten $v_1, \ldots, v_n$ und $n-1$
  Kanten $\{v_1,v_2\}, \{v_2,v_3\}, \ldots, \{v_{n-1},v_n\}$.
  Der \emph{Kreisgraph} $C_n$ entsteht aus $P_n$ durch Hinzuf\"ugen der
  Kante $\{v_n, v_1\}$; er hat $n$ Knoten und $n$ Kanten.
  Das Kantenma\ss{} betr\"agt $\kappa(P_n) = \frac{n}{n-1} \approx 1$ und
  $\kappa(C_n) = 1$.
\end{definition}

\begin{definition}[Baum, Wald]
  Ein \emph{Baum} ist ein zusammenh\"angender, kreisfreier Graph. Ein Baum
  mit $n$ Knoten hat genau $n-1$ Kanten. Ein \emph{Wald} ist ein kreisfreier
  (nicht notwendig zusammenh\"angender) Graph; seine Zusammenhangskomponenten
  sind B\"aume.
\end{definition}

\begin{definition}[Bipartiter Graph, vollst\"andiger bipartiter Graph]
  Ein Graph $G = (V,E)$ hei\ss{}t \emph{bipartit}, wenn die Knotenmenge $V$
  in zwei disjunkte Mengen $A$ und $B$ (die \emph{Bipartitionsklassen})
  zerlegt werden kann ($V = A \mathbin{\dot{\cup}} B$), so dass jede Kante
  einen Endknoten in $A$ und einen in $B$ hat, d.\,h.
  $E \subseteq \{\{a,b\} \mid a \in A, b \in B\}$.
  Der \emph{vollst\"andige bipartite Graph} $K_{r,s}$ mit $|A| = r$ und
  $|B| = s$ enth\"alt alle $r \cdot s$ m\"oglichen Kanten zwischen $A$ und $B$.
\end{definition}

\begin{satz}[Charakterisierung bipartiter Graphen]
  Ein Graph $G$ ist genau dann bipartit, wenn er keinen Kreis ungerader
  L\"ange enth\"alt.
\end{satz}

\begin{proof}
  ($\Rightarrow$) Sei $G$ bipartit mit Bipartitionsklassen $A$ und $B$.
  Jeder Kreis wechselt mit jeder Kante zwischen $A$ und $B$. Nach einer
  geraden Anzahl von Kanten ist man wieder in der Ausgangsklasse. Also hat
  jeder Kreis gerade L\"ange.
  ($\Leftarrow$) Sei $G$ zusammenh\"angend (der allgemeine Fall folgt
  komponentenweise). W\"ahle einen beliebigen Startknoten $s$ und f\"arbe
  alle Knoten gerade Distanz von $s$ mit Farbe $A$, alle mit ungerader
  Distanz mit Farbe $B$. Wenn eine Kante $\{u,v\}$ beide Endknoten in
  derselben Farbe h\"atte, so g\"abe es einen Weg von $s$ nach $u$, dann
  die Kante $\{u,v\}$, dann einen Weg von $v$ nach $s$ -- das ergebe einen
  Kreis ungerader L\"ange, Widerspruch zur Voraussetzung.
\end{proof}

\begin{definition}[Regul\"arer Graph, Petersen-Graph]
  Ein Graph hei\ss{}t \emph{$k$-regul\"ar}, wenn alle Knoten denselben Grad $k$
  besitzen. Bekannte Beispiele:
  \begin{itemize}
    \item Der \emph{Petersen-Graph} ist 3-regul\"ar mit 10 Knoten und 15
          Kanten; er ist ein wichtiges Gegenbeispiel in der Graphentheorie.
    \item Der \emph{Hyperkubus} $Q_d$ ist $d$-regul\"ar mit $2^d$ Knoten und
          $d \cdot 2^{d-1}$ Kanten; seine Knoten sind Bin\"arw\"orter der
          L\"ange $d$, zwei Knoten sind verbunden genau dann, wenn sie sich in
          genau einer Stelle unterscheiden.
  \end{itemize}
\end{definition}

\begin{definition}[Sterngraph $S_n$, Radgraph $W_n$]
  Der \emph{Sterngraph} $S_n = K_{1,n}$ besteht aus einem Zentrumsknoten,
  der mit $n$ Blattknoten verbunden ist; er hat $n+1$ Knoten und $n$ Kanten.
  Der \emph{Radgraph} $W_n$ entsteht aus $C_n$ durch Hinzuf\"ugen eines
  Zentrumsknotens, der mit allen $n$ Kreisknoten verbunden ist.
\end{definition}

\section{Isomorphie von Graphen}

\begin{definition}[Graphisomorphismus]
  Zwei Graphen $G_1 = (V_1, E_1)$ und $G_2 = (V_2, E_2)$ hei\ss{}en
  \emph{isomorph}, geschrieben $G_1 \cong G_2$, wenn es eine Bijektion
  $\phi: V_1 \to V_2$ gibt, so dass f\"ur alle $u, v \in V_1$ gilt:
  \[
    \{u,v\} \in E_1 \;\Leftrightarrow\; \{\phi(u), \phi(v)\} \in E_2.
  \]
  Die Abbildung $\phi$ hei\ss{}t \emph{Isomorphismus} von $G_1$ nach $G_2$.
\end{definition}

\begin{lemma}[Notwendige Bedingungen f\"ur Isomorphie]
  Sind $G_1 \cong G_2$, so gilt notwendigerweise:
  \begin{itemize}
    \item $|V(G_1)| = |V(G_2)|$ (gleiche Ordnung),
    \item $|E(G_1)| = |E(G_2)|$ (gleiche Gr\"o\ss{}e),
    \item Gleiche Gradfolge,
    \item Gleiche Anzahl von Zusammenhangskomponenten,
    \item Gleiche Taille und gleicher Umfang.
  \end{itemize}
  Diese Bedingungen sind jedoch \emph{nicht hinreichend} f\"ur Isomorphie:
  Es gibt nicht-isomorphe Graphen, die dieselbe Gradfolge besitzen.
\end{lemma}

\begin{beispiel}[Nicht-isomorphe Graphen mit gleicher Gradfolge]
  Die Graphen $C_4$ (Viereck) und $K_{1,3}$ plus isolierter Kante haben
  beide die Gradfolge $(2,2,2,2)$ bzw.\ $(3,1,1,1)$. Als ein subtileres
  Beispiel: Es gibt zwei nicht-isomorphe Graphen auf 6 Knoten mit Gradfolge
  $(3,3,2,2,2,2)$ -- n\"amlich den Kreis $C_6$ mit einer zus\"atzlichen
  Diagonale (zwei verschiedene Strukturen).
\end{beispiel}

\section{Aufgaben}

\begin{aufgabe}
  Zeigen Sie: In einem schlichten Graphen mit $n \geq 2$ Knoten gibt es
  stets zwei Knoten mit gleichem Grad.
  (\emph{Hinweis}: Schubfachprinzip -- die m\"oglichen Grade liegen in
  $\{0, 1, \ldots, n-1\}$, aber Grad $0$ und Grad $n-1$ k\"onnen nicht
  gleichzeitig auftreten.)
\end{aufgabe}

\begin{aufgabe}
  Bestimmen Sie alle nicht-isomorphen schlichten Graphen mit 4 Knoten und
  3 Kanten. (\emph{Hinweis}: Es gibt genau zwei.)
\end{aufgabe}

\begin{aufgabe}
  Berechnen Sie das Kantenma\ss{} $\kappa = |V|/|E|$ f\"ur $K_n$, $P_n$,
  $C_n$ und $K_{r,s}$. Welche Klasse hat das gr\"o\ss{}te bzw.\ kleinste
  Kantenma\ss{} asymptotisch?
\end{aufgabe}

\begin{aufgabe}
  Zeigen Sie, dass f\"ur jeden $k$-regul\"aren Graphen $G$ mit $n$ Knoten
  und $k$ ungerade gilt: $n$ ist gerade.
\end{aufgabe}

\begin{aufgabe}
  Sei $G$ ein Graph mit Gradfolge $(d_1, \ldots, d_n)$. Zeigen Sie, dass
  der Komplementgraph $\overline{G}$ die Gradfolge
  $(n-1-d_1, \ldots, n-1-d_n)$ hat.
\end{aufgabe}
\chapter{Graphen und Matrizen}
\label{chap:matrizen}

Die algebraische Behandlung von Graphen mittels Matrizen ist ein zentrales
Werkzeug sowohl f\"ur theoretische Analysen als auch f\"ur algorithmische
Berechnungen. Drei Matrizen sind besonders bedeutsam: die
\emph{Adjazenzmatrix}, die die Verbindungsstruktur kodiert; die
\emph{Inzidenzmatrix}, die die Beziehung zwischen Knoten und Kanten
beschreibt; sowie die \emph{Laplace-Matrix}, die spektrale Eigenschaften
des Graphen kodiert und mit der Anzahl der Spannb\"aume zusammenh\"angt.

\section{Die Adjazenzmatrix}

\begin{definition}[Adjazenzmatrix]
  F\"ur einen Graphen $G = (V,E)$ mit $n$ Knoten, deren Knoten wir mit
  $1, 2, \ldots, n$ benennen, ist die \emph{Adjazenzmatrix}
  $A \in \{0,1\}^{n \times n}$ definiert durch:
  \[
    A_{ij} = \begin{cases}
               1 & \text{falls } \{i,j\} \in E,\\
               0 & \text{sonst.}
             \end{cases}
  \]
  F\"ur ungerichtete Graphen ist $A$ symmetrisch ($A = A^\top$), da
  $\{i,j\} \in E \Leftrightarrow \{j,i\} \in E$.
\end{definition}

\begin{bemerkung}
  Die Adjazenzmatrix h\"angt von der Nummerierung der Knoten ab. Zwei
  Nummerierungen desselben Graphen liefern \"ahnliche Matrizen: Ist $A_1$
  die Adjazenzmatrix bez\"uglich Nummerierung~1 und $A_2$ bez\"uglich
  Nummerierung~2, so gilt $A_2 = P A_1 P^\top$ f\"ur eine geeignete
  Permutationsmatrix~$P$. Deshalb sind zwei Graphen genau dann isomorph,
  wenn ihre Adjazenzmatrizen durch eine Permutationsmatrix \"ahnlich sind.
\end{bemerkung}

\begin{beispiel}[Adjazenzmatrix des Pfadgraphen $P_4$]
  F\"ur $P_4$ mit Knoten $\{1,2,3,4\}$ und Kanten
  $\{\{1,2\},\{2,3\},\{3,4\}\}$ ist:
  \[
    A(P_4) = \begin{pmatrix}
               0 & 1 & 0 & 0\\
               1 & 0 & 1 & 0\\
               0 & 1 & 0 & 1\\
               0 & 0 & 1 & 0
             \end{pmatrix}.
  \]
  Die Zeilen- und Spaltensummen ergeben die Knotengrade:
  $\deg(1) = 1, \deg(2) = 2, \deg(3) = 2, \deg(4) = 1$.
\end{beispiel}

\subsection{Potenzen der Adjazenzmatrix und Boolean-Darstellung}

Das entscheidende Werkzeug f\"ur die sp\"atere Profilberechnung ist die
Verbindung zwischen Matrixpotenzen und Wegen im Graphen. Wir f\"uhren
zun\"achst die Boolean Matrixmultiplikation ein, die anstelle der
\"ublichen arithmetischen Operationen logische Operationen verwendet.

\begin{definition}[Boolean Matrixmultiplikation]
  F\"ur zwei Boolean-Matrizen $A, B \in \{0,1\}^{n \times n}$ ist die
  \emph{Boolean Matrixmultiplikation} definiert durch:
  \[
    C_{ij} = \bigvee_{k=1}^{n} (A_{ik} \wedge B_{kj}),
  \]
  wobei $\vee$ das logische ODER und $\wedge$ das logische UND bezeichnet.
  \"Aquivalent gilt: $C_{ij} = 1$ genau dann, wenn es ein $k$ mit
  $A_{ik} = 1$ und $B_{kj} = 1$ gibt.
\end{definition}

\begin{lemma}[Wege und Matrixpotenzen]
  \label{lem:wege-matrizen}
  Sei $A$ die Adjazenzmatrix eines Graphen $G = (V,E)$ und $A^k$ die
  $k$-te Potenz unter Boolean Matrixmultiplikation. Dann gilt:
  \[
    (A^k)_{ij} = 1 \;\Leftrightarrow\;
    \text{es existiert ein Weg der L\"ange } k \text{ von } i \text{ nach } j.
  \]
\end{lemma}

\begin{proof}
  Vollst\"andige Induktion \"uber $k \geq 1$.

  \textbf{Induktionsanfang} ($k = 1$): Es gilt $(A^1)_{ij} = A_{ij} = 1$
  genau dann, wenn $\{i,j\} \in E$, d.\,h. ein Weg der L\"ange~1 von $i$
  nach $j$ existiert. \checkmark

  \textbf{Induktionsschritt} ($k \to k+1$): Sei die Aussage f\"ur $k$
  bewiesen. Dann gilt:
  \[
    (A^{k+1})_{ij} = \bigvee_{l=1}^{n}\bigl((A^k)_{il} \wedge A_{lj}\bigr).
  \]
  Dieses Ergebnis ist genau dann $= 1$, wenn es ein $l$ gibt mit
  $(A^k)_{il} = 1$ und $A_{lj} = 1$. Nach Induktionsvoraussetzung bedeutet
  $(A^k)_{il} = 1$, dass ein Weg der L\"ange $k$ von $i$ nach $l$ existiert.
  Und $A_{lj} = 1$ bedeutet, dass die Kante $\{l,j\}$ existiert. Zusammen
  ergibt das einen Weg der L\"ange $k+1$ von $i$ nach $j$.
  Umgekehrt: Existiert ein Weg $(i = v_0, v_1, \ldots, v_k, v_{k+1} = j)$
  der L\"ange $k+1$, so ist $(v_0, \ldots, v_k)$ ein Weg der L\"ange $k$ von
  $i$ nach $v_k$ (also $(A^k)_{i,v_k} = 1$ nach Induktionsvoraussetzung) und
  $\{v_k, j\} \in E$ (also $A_{v_k,j} = 1$). Damit ist
  $(A^{k+1})_{ij} \geq (A^k)_{i,v_k} \wedge A_{v_k,j} = 1$.
\end{proof}

\begin{korollar}
  F\"ur die Anzahl der Wege der L\"ange $k$ zwischen $i$ und $j$ gilt bei
  der gew\"ohnlichen (nicht-Booleschen) Matrixmultiplikation:
  $(A^k)_{ij}$ ist gleich der Anzahl der (nicht notwendig einfachen)
  Wege der L\"ange $k$ von $i$ nach $j$.
\end{korollar}

\begin{satz}[Erreichbarkeitsmatrix]
  \label{satz:erreichbarkeit}
  Die \emph{Erreichbarkeitsmatrix} $R \in \{0,1\}^{n \times n}$ mit
  $R_{ij} = 1$ falls ein Weg von $i$ nach $j$ existiert (sonst $0$)
  erf\"ullt:
  \[
    R = A \vee A^2 \vee A^3 \vee \cdots \vee A^{n-1},
  \]
  wobei $\vee$ komponentenweise das logische ODER bezeichnet.
\end{satz}

\begin{proof}
  Ein Weg von $i$ nach $j$ existiert genau dann, wenn ein Weg einer
  L\"ange $k \in \{1, \ldots, n-1\}$ von $i$ nach $j$ existiert. Die
  Schranke $n-1$ ergibt sich, weil jeder einfache Weg in einem Graphen
  mit $n$ Knoten h\"ochstens $n-1$ Kanten hat (mehr Kanten w\"urde
  bedeuten, dass ein Knoten doppelt besucht wird). Nach
  Lemma~\ref{lem:wege-matrizen} ist $(A^k)_{ij} = 1$ \"aquivalent zum
  Existieren eines Weges der L\"ange $k$. Die komponentenweise ODER-Verkn\"upfung
  liefert genau dann $1$, wenn mindestens eine der Potenzen $1$ ist.
\end{proof}

\subsection{Zerlegbare Matrizen und Zusammenhang}

\begin{definition}[Zerlegbare und irreduzible Matrix]
  Eine quadratische Matrix $A \in \mathbb{R}^{n \times n}$ hei\ss{}t
  \emph{zerlegbar}, wenn es eine Permutationsmatrix $P$ gibt, sodass
  $P A P^\top$ Blockdreiecksgestalt besitzt:
  \[
    P A P^\top = \begin{pmatrix} B & C \\ 0 & D \end{pmatrix},
  \]
  wobei $B$ und $D$ quadratische Matrizen sind. Andernfalls hei\ss{}t $A$
  \emph{irreduzibel}.
\end{definition}

\begin{satz}
  Die Adjazenzmatrix $A$ eines Graphen $G$ ist genau dann irreduzibel,
  wenn $G$ zusammenh\"angend ist.
\end{satz}

\section{Inzidenzmatrix und Gradmatrix}

\begin{definition}[Inzidenzmatrix]
  F\"ur einen Graphen $G = (V,E)$ mit $n$ Knoten und $m$ Kanten ist die
  \emph{Inzidenzmatrix} $B \in \{0,1\}^{n \times m}$ definiert durch:
  \[
    B_{ie} = \begin{cases}
               1 & \text{falls Knoten } i \text{ Endknoten der Kante } e \text{ ist},\\
               0 & \text{sonst.}
             \end{cases}
  \]
  Jede Spalte von $B$ hat genau zwei Einsen (die beiden Endknoten der
  zugeh\"origen Kante). Es gilt $B \mathbf{1} = \deg$ (Gradvektor).
\end{definition}

\begin{definition}[Gradmatrix]
  Die \emph{Gradmatrix} $D_G$ ist die $n \times n$-Diagonalmatrix mit
  $(D_G)_{ii} = \deg(i)$ und $(D_G)_{ij} = 0$ f\"ur $i \neq j$.
\end{definition}

\section{Abst\"ande in Graphen und das Kantenma\ss{}}

\begin{definition}[Abstand, Distanzmatrix]
  Der \emph{Abstand} $d(u,v)$ zweier Knoten $u, v \in V$ ist die L\"ange
  eines k\"urzesten Weges von $u$ nach $v$ (bzw. $\infty$, falls kein Weg
  existiert; $d(v,v) = 0$). Die \emph{Distanzmatrix}
  $D \in (\mathbb{N}_0 \cup \{\infty\})^{n \times n}$ hat den Eintrag
  $D_{ij} = d(i,j)$.
\end{definition}

\begin{definition}[Durchmesser, Radius, Exzentrizit\"at, Zentrum]
  Die \emph{Exzentrizit\"at} eines Knotens $v$ ist
  $\varepsilon(v) = \max_{u \in V} d(u,v)$.
  Der \emph{Durchmesser} ist $\mathrm{diam}(G) = \max_{v \in V} \varepsilon(v)
  = \max_{u,v \in V} d(u,v)$.
  Der \emph{Radius} ist $\mathrm{rad}(G) = \min_{v \in V} \varepsilon(v)$.
  Das \emph{Zentrum} ist die Menge der Knoten minimaler Exzentrizit\"at:
  $\mathrm{Z}(G) = \{v \in V \mid \varepsilon(v) = \mathrm{rad}(G)\}$.
\end{definition}

\begin{satz}
  F\"ur jeden zusammenh\"angenden Graphen $G$ gilt:
  $\mathrm{rad}(G) \leq \mathrm{diam}(G) \leq 2\,\mathrm{rad}(G)$.
\end{satz}

\begin{proof}
  Die erste Ungleichung $\mathrm{rad}(G) \leq \mathrm{diam}(G)$ folgt
  unmittelbar aus den Definitionen: Der Radius ist das Minimum, der
  Durchmesser das Maximum der Exzentrizit\"aten.

  F\"ur die zweite Ungleichung sei $c \in \mathrm{Z}(G)$ ein Zentrumsknoten
  mit $\varepsilon(c) = \mathrm{rad}(G)$. F\"ur beliebige $u, v \in V$ gilt:
  \[
    d(u,v) \leq d(u,c) + d(c,v) \leq \varepsilon(c) + \varepsilon(c)
            = 2\,\mathrm{rad}(G). \qedhere
  \]
\end{proof}

\begin{definition}[Kantenma\ss{}]
  \label{def:kantenmass}
  Das \emph{Kantenma\ss{}} $\kappa$ eines Graphen $G = (V,E)$ ist:
  \[
    \kappa = \frac{|V|}{|E|} = \frac{n}{m}.
  \]
  Ein kleines $\kappa$ bedeutet viele Verbindungen pro Knoten (dichter Graph);
  ein gro\ss{}es $\kappa$ kennzeichnet sp\"arlich verbundene Topologien. Das
  Kantenma\ss{} ist eine der drei Komponenten des Graphprofils
  (vgl.~Kapitel~\ref{chap:algorithmen}).
\end{definition}

\begin{lemma}[Extremwerte des Kantenma\ss{}es]
  \label{lem:kappa-extrem}
  F\"ur spezielle Graphklassen gilt:
  \begin{itemize}
    \item Vollst\"andiger Graph $K_n$:
          $\kappa(K_n) = \frac{2}{n-1} \xrightarrow{n\to\infty} 0$
    \item Pfadgraph $P_n$:
          $\kappa(P_n) = \frac{n}{n-1} \xrightarrow{n\to\infty} 1$
    \item Sterngraph $S_n = K_{1,n}$:
          $\kappa(S_n) = \frac{n+1}{n} \xrightarrow{n\to\infty} 1$
    \item Kreisgraph $C_n$: $\kappa(C_n) = 1$ f\"ur alle $n \geq 3$
    \item Vollst\"andiger bipartiter Graph $K_{n,n}$:
          $\kappa(K_{n,n}) = \frac{2}{n} \xrightarrow{n\to\infty} 0$
    \item Graph ohne Kanten: $\kappa = \infty$
  \end{itemize}
\end{lemma}

\section{Spannb\"aume und die Laplace-Matrix}

\begin{definition}[Spannbaum]
  Ein \emph{Spannbaum} eines zusammenh\"angenden Graphen $G$ ist ein
  Teilgraph $T$, der alle Knoten von $G$ enth\"alt und ein Baum ist.
  Jeder zusammenh\"angende Graph besitzt mindestens einen Spannbaum.
\end{definition}

\begin{definition}[Laplace-Matrix]
  Die \emph{Laplace-Matrix} (auch \emph{Kirchhoff-Matrix}) eines Graphen
  $G$ ist definiert als $L_G = D_G - A$, wobei $D_G$ die Gradmatrix und
  $A$ die Adjazenzmatrix ist. Es gilt:
  \[
    (L_G)_{ij} = \begin{cases}
                   \deg(i) & \text{falls } i = j,\\
                   -1      & \text{falls } \{i,j\} \in E,\\
                   0       & \text{sonst.}
                 \end{cases}
  \]
  Die Laplace-Matrix erf\"ullt $L_G = B B^\top$ f\"ur die Inzidenzmatrix $B$
  (bei geeigneter Orientierung der Kanten).
\end{definition}

\begin{satz}[Kirchhoff, 1847 -- Satz vom Matrix-Baum]
  Sei $G$ ein zusammenh\"angender Graph mit Laplace-Matrix $L_G$. Dann ist
  die Anzahl der Spannb\"aume von $G$ gleich jedem Kofaktor von $L_G$,
  d.\,h. gleich dem $(i,i)$-Kofaktor $\det(L_G^{(i)})$, wobei $L_G^{(i)}$
  die Matrix $L_G$ ohne die $i$-te Zeile und $i$-te Spalte ist (f\"ur
  beliebiges $i$).
\end{satz}

\begin{beispiel}[Anzahl der Spannb\"aume von $K_4$]
  F\"ur den vollst\"andigen Graphen $K_4$ betr\"agt die Laplace-Matrix:
  \[
    L_{K_4} = \begin{pmatrix}
                3 & -1 & -1 & -1\\
               -1 &  3 & -1 & -1\\
               -1 & -1 &  3 & -1\\
               -1 & -1 & -1 &  3
              \end{pmatrix}.
  \]
  Der Kofaktor $\det(L_{K_4}^{(1)}) = 16 = 4^{4-2}$. Nach der
  Cayleyschen Formel hat $K_n$ genau $n^{n-2}$ Spannb\"aume; f\"ur
  $n = 4$ sind das $4^2 = 16$. \checkmark
\end{beispiel}

\begin{satz}[Eigenschaften der Laplace-Matrix]
  F\"ur jeden Graphen $G$ gilt:
  \begin{itemize}
    \item $L_G$ ist symmetrisch und positiv semidefinit.
    \item Alle Eigenwerte von $L_G$ sind reell und nicht-negativ:
          $0 = \lambda_1 \leq \lambda_2 \leq \cdots \leq \lambda_n$.
    \item Die algebraische Vielfachheit des Eigenwerts $0$ ist gleich der
          Anzahl der Zusammenhangskomponenten von $G$.
    \item Der kleinste positive Eigenwert $\lambda_2$ (der
          \emph{Fiedler-Wert}) ist ein Ma\ss{} f\"ur die Konnektivit\"at
          des Graphen.
  \end{itemize}
\end{satz}

\section{Aufgaben}

\begin{aufgabe}
  Berechnen Sie $A^2$ und $A^3$ f\"ur den Pfadgraphen $P_4$ und
  interpretieren Sie jeden Eintrag $(A^k)_{ij}$ graphentheoretisch.
\end{aufgabe}

\begin{aufgabe}
  Bestimmen Sie Durchmesser, Radius und Zentrum f\"ur $K_n$, $P_n$,
  $C_n$ und $K_{r,s}$.
\end{aufgabe}

\begin{aufgabe}
  Berechnen Sie die Laplace-Matrix von $P_4$ und $C_4$.
  Verifizieren Sie, dass die Spaltensummen null sind.
\end{aufgabe}

\begin{aufgabe}
  Zeigen Sie: $\mathrm{rad}(G) \leq \mathrm{diam}(G) \leq 2\,\mathrm{rad}(G)$
  und geben Sie f\"ur jede Seite ein Beispiel an, bei dem Gleichheit gilt.
\end{aufgabe}

\begin{aufgabe}
  Bestimmen Sie die Erreichbarkeitsmatrix $R$ des Graphen mit
  $V = \{1,2,3,4\}$ und $E = \{\{1,2\},\{2,3\},\{3,4\}\}$ durch
  sukzessive Berechnung von $A, A^2, A^3$.
\end{aufgabe}
\chapter{Planare Graphen und die Eulersche Polyederformel}

Ein Graph hei\ss{}t \emph{planar}, wenn er in der Ebene so gezeichnet werden
kann, dass sich keine zwei Kanten \"uberschneiden. Diese Eigenschaft ist
fundamental f\"ur Schaltkreisdesign, Kartographie und Computergrafik. Das
wichtigste Resultat dieses Kapitels -- die Eulersche Polyederformel -- stellt
eine elegante Beziehung zwischen Knoten, Kanten und Fl\"achen her und erlaubt
es, viele planare Graphen durch Kantenanzahlschranken zu charakterisieren.

\section{Planare Einbettungen}

\begin{definition}[Planarer Graph, planare Einbettung, Fl\"achen]
  Ein Graph $G$ hei\ss{}t \emph{planar}, wenn er eine \emph{planare
  Einbettung} besitzt: eine Darstellung in der Ebene $\mathbb{R}^2$, bei
  der die Knoten als Punkte und die Kanten als Jordan-Kurven (stetige
  injektive Kurven) dargestellt werden, sodass sich zwei Kanten nur in
  gemeinsamen Endknoten schneiden.

  Eine planare Einbettung zerlegt die Ebene in zusammenh\"angende offene
  Gebiete, die \emph{Fl\"achen} (Facetten) hei\ss{}en. Genau eine dieser
  Fl\"achen ist unbeschr\"ankt; sie hei\ss{}t \emph{\"au\ss{}ere Fl\"ache}.
\end{definition}

\begin{bemerkung}
  Nicht jede geradlinige Zeichnung eines planaren Graphen ist planar, aber
  der Satz von Fary (1948) besagt, dass jeder planare Graph eine planare
  Einbettung besitzt, bei der alle Kanten geradlinige Segmente sind.
\end{bemerkung}

\begin{definition}[Kreuzungszahl]
  Die \emph{Kreuzungszahl} $\mathrm{cr}(G)$ eines Graphen ist die minimale
  Anzahl von Kantenkreuzungen \"uber alle m\"oglichen Zeichnungen von $G$ in
  der Ebene. Es gilt $\mathrm{cr}(G) = 0$ genau dann, wenn $G$ planar ist.
\end{definition}

\section{Die Eulersche Polyederformel}

\begin{satz}[Eulersche Polyederformel, Euler 1758]
  \label{satz:euler}
  F\"ur jeden zusammenh\"angenden planaren Graphen $G$ mit $n$ Knoten,
  $m$ Kanten und $f$ Fl\"achen gilt:
  \[
    n - m + f = 2.
  \]
\end{satz}

\begin{proof}
  Wir beweisen die Formel durch Induktion \"uber die Anzahl der Kanten $m$.

  \textbf{Induktionsanfang}: Hat $G$ einen einzigen Knoten und keine Kanten
  ($m = 0$), so gibt es $n = 1$ Knoten und $f = 1$ Fl\"ache (die gesamte
  Ebene). Also $1 - 0 + 1 = 2$. \checkmark

  \textbf{Induktionsschritt}: Sei $m \geq 1$ und die Formel f\"ur alle
  zusammenh\"angenden planaren Graphen mit weniger Kanten bewiesen.
  \begin{itemize}
    \item Falls $G$ eine Br\"ucke $e$ enth\"alt: Dann trennt $G - e$ in zwei
          Komponenten. Aber da $G$ zusammenh\"angend war, hat $G - e$ genau
          zwei Zusammenhangskomponenten, die beide zusammenh\"angend sind.
          Tats\"achlich hat $G - e$ weniger Kanten, ist aber nicht mehr
          zusammenh\"angend. Daher betrachten wir stattdessen den Fall:
    \item Falls $G$ keinen Kreis enth\"alt: Dann ist $G$ ein Baum mit
          $m = n-1$ Kanten und $f = 1$ Fl\"ache. Dann $n - (n-1) + 1 = 2$.
          \checkmark
    \item Falls $G$ einen Kreis enth\"alt: Sei $e$ eine Kante dieses Kreises.
          Das Entfernen von $e$ ergibt $G' = G - e$. $G'$ ist immer noch
          zusammenh\"angend (da $e$ in einem Kreis lag, also keine Br\"ucke
          war). $G'$ hat $n$ Knoten, $m-1$ Kanten und $f-1$ Fl\"achen (das
          Entfernen von $e$ vereinigt zwei benachbarte Fl\"achen zu einer).
          Nach Induktionsvoraussetzung: $n - (m-1) + (f-1) = 2$, also
          $n - m + f = 2$. \checkmark \qedhere
  \end{itemize}
\end{proof}

\begin{korollar}[Kantenschranken f\"ur planare Graphen]
  \label{kor:kantenschranken}
  \begin{itemize}
    \item F\"ur jeden einfachen zusammenh\"angenden planaren Graphen mit
          $n \geq 3$ Knoten gilt: $m \leq 3n - 6$.
    \item Falls $G$ zus\"atzlich bipartit ist ($n \geq 3$):
          $m \leq 2n - 4$.
    \item Aus $m \leq 3n-6$ folgt: $\kappa = n/m \geq n/(3n-6) \approx 1/3$
          f\"ur planare Graphen.
  \end{itemize}
\end{korollar}

\begin{proof}
  Jede Fl\"ache ist von mindestens 3 Kanten begrenzt. Da jede Kante zu
  h\"ochstens 2 Fl\"achen geh\"ort, gilt $3f \leq 2m$. Mit der Eulerschen
  Formel $f = 2 - n + m$ ergibt sich $3(2-n+m) \leq 2m$, also $m \leq 3n-6$.
  F\"ur bipartite Graphen (Taille $\geq 4$) gilt $4f \leq 2m$, was zu
  $m \leq 2n-4$ f\"uhrt.
\end{proof}

\section{Nichtplanare Graphen und der Satz von Kuratowski}

\begin{satz}[$K_5$ und $K_{3,3}$ sind nicht planar]
  Die Graphen $K_5$ und $K_{3,3}$ sind nicht planar.
\end{satz}

\begin{proof}
  $K_5$: Hat $n = 5$, $m = 10$. W\"are $K_5$ planar, so m\"usste nach
  Korollar~\ref{kor:kantenschranken} gelten: $m \leq 3n - 6 = 9$.
  Da $10 > 9$, ist $K_5$ nicht planar.

  $K_{3,3}$: Hat $n = 6$, $m = 9$. Da $K_{3,3}$ bipartit ist, m\"usste
  gelten: $m \leq 2n - 4 = 8$. Da $9 > 8$, ist $K_{3,3}$ nicht planar.
\end{proof}

\begin{definition}[Unterteilung]
  Eine \emph{Unterteilung} $G'$ eines Graphen $G$ entsteht, indem man
  Kanten von $G$ durch Pfade ersetzt (d.\,h. man f\"ugt auf einer Kante
  zus\"atzliche Knoten vom Grad~2 ein).
\end{definition}

\begin{satz}[Kuratowski, 1930]
  Ein Graph ist genau dann planar, wenn er keinen Teilgraphen enth\"alt,
  der eine Unterteilung von $K_5$ oder $K_{3,3}$ ist.
\end{satz}

\begin{satz}[Wagner, 1937]
  Ein Graph ist genau dann planar, wenn er weder $K_5$ noch $K_{3,3}$ als
  Minor enth\"alt (ein Graph $H$ ist ein \emph{Minor} von $G$, wenn $H$ durch
  Kantenkontraktionen aus einem Teilgraphen von $G$ entsteht).
\end{satz}

\section{Der duale Graph}

\begin{definition}[Dualer Graph]
  Sei $G$ ein zusammenh\"angender planarer Graph mit einer planaren Einbettung.
  Der \emph{duale Graph} $G^*$ hat als Knotenmenge die Fl\"achen von $G$;
  f\"ur jede Kante $e \in E(G)$ gibt es eine Kante $e^*$ in $G^*$, die die
  beiden Fl\"achen verbindet, die $e$ trennt (bei einer Br\"ucke wird
  dieselbe Fl\"ache zweifach verbunden, was eine Schlinge in $G^*$ ergibt).
\end{definition}

\begin{satz}
  F\"ur den dualen Graphen $G^*$ eines zusammenh\"angenden planaren Graphen
  $G$ gilt $(G^*)^* \cong G$.
\end{satz}

\section{Aufgaben}

\begin{aufgabe}
  \"Uberpr\"ufen Sie die Eulersche Formel f\"ur den W\"urfelgraphen $Q_3$
  (8 Knoten, 12 Kanten) und bestimmen Sie die Anzahl seiner Fl\"achen.
\end{aufgabe}

\begin{aufgabe}
  Zeigen Sie, dass der Petersen-Graph nicht planar ist.
  (\emph{Hinweis}: Zeigen Sie, dass er $K_{3,3}$ als Minor enth\"alt.)
\end{aufgabe}

\begin{aufgabe}
  Zeigen Sie: Jeder planare Graph enth\"alt mindestens einen Knoten mit
  Grad $\leq 5$.
\end{aufgabe}

%=======================================================
\chapter{Unabh\"angige Mengen, Cliquen und Matchings}
%=======================================================

Unabh\"angige Mengen und Matchings sind grundlegende Konzepte f\"ur
Optimierungsprobleme in Graphen. Sie modellieren Konfliktsituationen
(keine zwei ausgew\"ahlten Elemente d\"urfen benachbart sein) und
Zuordnungsprobleme (maximale Paarungen zwischen Knoten). Anwendungen
reichen von der Ressourcenzuteilung \"uber Stundenplanerstellung bis hin
zur DNA-Sequenzierung.

\section{Unabh\"angige Knotenmengen}

\begin{definition}[Unabh\"angige Knotenmenge, Unabh\"angigkeitszahl]
  Eine Teilmenge $S \subseteq V$ hei\ss{}t \emph{unabh\"angige Knotenmenge}
  (oder \emph{stabile Menge}), wenn keine zwei Knoten in $S$ durch eine
  Kante verbunden sind:
  $\forall u, v \in S: \{u,v\} \notin E.$
  Die \emph{Unabh\"angigkeitszahl} $\alpha(G)$ ist die maximale Gr\"o\ss{}e
  einer unabh\"angigen Knotenmenge.
\end{definition}

\begin{definition}[Clique, Cliquenzahl]
  Eine Teilmenge $C \subseteq V$ hei\ss{}t \emph{Clique}, wenn je zwei
  Knoten in $C$ durch eine Kante verbunden sind: $C$ induziert einen
  vollst\"andigen Teilgraphen $K_{|C|}$. Die \emph{Cliquenzahl} $\omega(G)$
  ist die maximale Gr\"o\ss{}e einer Clique.
\end{definition}

\begin{satz}[Komplement-Beziehung]
  F\"ur jeden Graphen $G$ gilt $\alpha(G) = \omega(\overline{G})$ und
  $\omega(G) = \alpha(\overline{G})$.
\end{satz}

\begin{proof}
  Eine unabh\"angige Menge in $G$ ist genau eine Menge, bei der keine zwei
  Knoten eine Kante haben. Das bedeutet: alle Paare, die keine Kante in $G$
  bilden, bilden eine Kante in $\overline{G}$. Also ist eine unabh\"angige
  Menge in $G$ genau eine Clique in $\overline{G}$.
\end{proof}

\begin{definition}[Knotenbedeckung, Bedeckungszahl]
  Eine Teilmenge $C \subseteq V$ hei\ss{}t \emph{Knotenbedeckung}, wenn
  jede Kante $\{u,v\} \in E$ mindestens einen Endknoten in $C$ hat:
  $\forall \{u,v\} \in E: u \in C \text{ oder } v \in C.$
  Die \emph{Bedeckungszahl} $\tau(G)$ ist die minimale Gr\"o\ss{}e einer
  Knotenbedeckung.
\end{definition}

\begin{satz}[Gallai, 1959]
  F\"ur jeden Graphen $G$ ohne isolierte Knoten gilt:
  $\alpha(G) + \tau(G) = n = |V|.$
\end{satz}

\begin{proof}
  Sei $S \subseteq V$ eine unabh\"angige Menge. Dann ist $V \setminus S$
  eine Knotenbedeckung: Da $S$ unabh\"angig ist, hat jede Kante
  $\{u,v\} \in E$ mindestens einen Endknoten in $V \setminus S$.
  Umgekehrt: Ist $C$ eine Knotenbedeckung, so ist $V \setminus C$
  unabh\"angig. Also ist $S$ maximal unabh\"angig genau dann, wenn
  $V \setminus S$ minimal bedeckend ist. Damit gilt
  $\alpha(G) = n - \tau(G)$.
\end{proof}

\section{Matchings}

\begin{definition}[Matching, S\"attigungsgrad, Matchingzahl]
  Eine Teilmenge $M \subseteq E$ hei\ss{}t \emph{Matching}, wenn keine zwei
  Kanten in $M$ einen gemeinsamen Endknoten teilen. Ein Matching
  hei\ss{}t \emph{perfekt}, wenn jeder Knoten genau einem Endknoten einer
  Matchingkante ist (nur m\"oglich bei geradem $n$). Die
  \emph{Matchingzahl} $\nu(G)$ ist die Gr\"o\ss{}e eines Matchings
  maximaler Gr\"o\ss{}e (\emph{maximales Matching}).
\end{definition}

\begin{definition}[Augmentierender Weg]
  Sei $M$ ein Matching. Ein \emph{augmentierender Weg} bez\"uglich $M$ ist
  ein Weg, der an einem unges\"attigten Knoten beginnt und endet und dessen
  Kanten abwechselnd nicht zu $M$ und zu $M$ geh\"oren.
\end{definition}

\begin{satz}[Berge, 1957]
  Ein Matching $M$ ist genau dann ein Maximum-Matching (maximaler Gr\"o\ss{}e),
  wenn kein augmentierender Weg bez\"uglich $M$ existiert.
\end{satz}

\begin{proof}
  ($\Rightarrow$) Ist $M$ ein Maximum-Matching und $P$ ein augmentierender
  Weg, so kann man $M$ durch die Kanten von $P$ verbessern (symmetrische
  Differenz $M \triangle P$ ist ein gr\"o\ss{}eres Matching) -- Widerspruch.

  ($\Leftarrow$) Existiert kein augmentierender Weg, so sei $M^*$ ein
  beliebiges Maximum-Matching. Die symmetrische Differenz $M \triangle M^*$
  besteht aus Pfaden und Kreisen, bei denen Kanten von $M$ und $M^*$
  abwechseln. Da $|M^*| > |M|$ angenommen, gibt es eine Komponente mit
  mehr Kanten aus $M^*$ als aus $M$ -- das ist ein augmentierender Weg
  f\"ur $M$, Widerspruch.
\end{proof}

\begin{satz}[K\"onig, 1931]
  In jedem bipartiten Graphen gilt: $\nu(G) = \tau(G)$.
  (Die Gr\"o\ss{}e eines maximalen Matchings ist gleich der Gr\"o\ss{}e
  einer minimalen Knotenbedeckung.)
\end{satz}

\begin{satz}[Hall, 1935]
  Sei $G = (A \cup B, E)$ ein bipartiter Graph. Dann besitzt $G$ ein Matching,
  das alle Knoten von $A$ s\"attigt, genau dann, wenn f\"ur alle
  $S \subseteq A$ gilt: $|N(S)| \geq |S|$ (Heiratsbed ingung).
\end{satz}

\section{Aufgaben}

\begin{aufgabe}
  Bestimmen Sie $\alpha$, $\omega$, $\nu$ und $\tau$ f\"ur den
  Petersen-Graphen.
\end{aufgabe}

\begin{aufgabe}
  Zeigen Sie: In jedem $k$-regul\"aren bipartiten Graphen existiert ein
  perfektes Matching.
\end{aufgabe}

\begin{aufgabe}
  Geben Sie einen Graphen an, bei dem $\omega(G) < \chi(G)$.
  (\emph{Hinweis}: Ber\"ucksichtigen Sie den F\"unf-Kreis $C_5$.)
\end{aufgabe}

%=======================================================
\chapter{F\"arbungen von Graphen}
%=======================================================

Graphf\"arbungen modellieren Konfliktsituationen, in denen benachbarte
Elemente verschiedene Zust\"ande einnehmen m\"ussen: Kanalzuweisung im
Mobilfunk, Registerallokation in Compilern, Stundenplanerstellung,
Kartographierung. Das zentrale Konzept ist die chromatische Zahl -- die
minimale Anzahl von Farben, die ben\"otigt wird, um einen Graphen zul\"assig
zu f\"arben.

\section{Knotenf\"arbungen}

\begin{definition}[Zul\"assige Knotenf\"arbung, chromatische Zahl]
  Eine \emph{(zul\"assige) $k$-Knotenf\"arbung} von $G$ ist eine Funktion
  $c: V \to \{1, 2, \ldots, k\}$ (die $k$ Farben), sodass je zwei
  benachbarte Knoten verschiedene Farben erhalten:
  $\forall \{u,v\} \in E: c(u) \neq c(v).$
  Die \emph{chromatische Zahl} $\chi(G)$ ist das minimale $k$, f\"ur das
  eine zul\"assige $k$-F\"arbung von $G$ existiert.
\end{definition}

\begin{bemerkung}
  Offensichtlich gilt: $\chi(G) = 1$ genau dann, wenn $G$ keine Kanten hat.
  $\chi(G) = 2$ genau dann, wenn $G$ mindestens eine Kante und bipartit ist.
  $\chi(K_n) = n$ und $\chi(C_{2k}) = 2$, $\chi(C_{2k+1}) = 3$.
\end{bemerkung}

\begin{satz}[Schranken f\"ur die chromatische Zahl]
  F\"ur jeden Graphen $G$ gilt:
  \[
    \omega(G) \leq \chi(G) \leq \Delta(G) + 1.
  \]
\end{satz}

\begin{proof}
  \emph{Untere Schranke}: Jede Clique der Gr\"o\ss{}e $\omega$ ben\"otigt
  $\omega$ verschiedene Farben.

  \emph{Obere Schranke} (gieriger Algorithmus): Ordne die Knoten in
  beliebiger Reihenfolge. Weise jedem Knoten die kleinste Farbe zu, die
  nicht von einem bereits gef\"arbten Nachbarn verwendet wird. Da jeder
  Knoten h\"ochstens $\Delta(G)$ Nachbarn hat, wird h\"ochstens die
  $(\Delta(G)+1)$-te Farbe ben\"otigt.
\end{proof}

\begin{satz}[Brooks, 1941]
  F\"ur jeden zusammenh\"angenden Graphen $G$, der weder ein vollst\"andiger
  Graph noch ein ungerader Kreis ist, gilt:
  $\chi(G) \leq \Delta(G).$
\end{satz}

\begin{definition}[Kantef\"arbung, chromatischer Index]
  Eine \emph{$k$-Kantenf\"arbung} weist jeder Kante eine Farbe aus
  $\{1,\ldots,k\}$ zu, sodass keine zwei inzidenten Kanten dieselbe Farbe
  haben. Der \emph{chromatische Index} $\chi'(G)$ ist das minimale $k$.
\end{definition}

\begin{satz}[Vizing, 1964]
  F\"ur jeden schlichten Graphen gilt:
  $\Delta(G) \leq \chi'(G) \leq \Delta(G) + 1.$
\end{satz}

\section{F\"arbungen planarer Graphen}

\begin{satz}[F\"unf-Farben-Satz]
  Jeder planare Graph ist 5-f\"arbbar.
\end{satz}

\begin{proof}[Beweisidee]
  Induktion \"uber $n$. Jeder planare Graph enth\"alt einen Knoten $v$ mit
  Grad $\leq 5$ (folgt aus $m \leq 3n-6$). Nach Entfernen von $v$ ist der
  resultierende Graph nach Induktionsvoraussetzung 5-f\"arbbar. Falls
  $\deg(v) \leq 4$, wird eine freie Farbe f\"ur $v$ gefunden. Im Fall
  $\deg(v) = 5$ mit allen 5 Nachbarn in verschiedenen Farben verwendet man
  den Kempe-Ketten-Argument, um die F\"arbung umzustrukturieren.
\end{proof}

\begin{satz}[Vier-Farben-Satz, Appel \& Haken, 1976]
  Jeder planare Graph ist 4-f\"arbbar: $\chi(G) \leq 4$.
\end{satz}

\begin{bemerkung}
  Der Vier-Farben-Satz war das erste gro\ss{}e mathematische Theorem, das
  wesentlich mit Hilfe eines Computers bewiesen wurde. Appel und Haken
  reduzierten das Problem auf die \"Uberpr\"ufung von 1936 Konfigurationen
  mit einem Computer. Ein vereinfachter Beweis wurde 1997 von Robertson,
  Sanders, Seymour und Thomas gegeben.
\end{bemerkung}

\section{Das chromatische Polynom}

\begin{definition}[Chromatisches Polynom]
  Das \emph{chromatische Polynom} $P(G, k)$ z\"ahlt die Anzahl der
  zul\"assigen $k$-F\"arbungen von $G$ f\"ur $k \in \mathbb{N}$.
  Es ist ein Polynom in $k$ vom Grad $n = |V|$ mit ganzzahligen
  Koeffizienten und f\"uhrendem Term $k^n$.
\end{definition}

\begin{beispiel}[Chromatische Polynome]
  \begin{itemize}
    \item Leerer Graph ($m=0$): $P(G, k) = k^n$
    \item Vollst\"andiger Graph: $P(K_n, k) = k(k-1)(k-2)\cdots(k-n+1)$
    \item Baum mit $n$ Knoten: $P(T, k) = k(k-1)^{n-1}$
    \item Kreis $C_n$: $P(C_n, k) = (k-1)^n + (-1)^n(k-1)$
  \end{itemize}
\end{beispiel}

\begin{satz}[Dekompositionsgleichung]
  F\"ur jede Kante $e = \{u,v\} \in E$:
  \[
    P(G, k) = P(G - e, k) - P(G / e, k),
  \]
  wobei $G/e$ den durch Kontraktion von $e$ entstehenden Graphen bezeichnet.
\end{satz}

\begin{proof}
  Die zul\"assigen $k$-F\"arbungen von $G - e$ zerfallen in zwei Klassen:
  solche, bei denen $c(u) \neq c(v)$ (diese sind auch zul\"assig in $G$),
  und solche, bei denen $c(u) = c(v)$ (diese entsprechen zul\"assigen
  F\"arbungen von $G/e$, bei dem $u$ und $v$ zu einem Knoten verschmolzen
  sind). Daher: $P(G-e, k) = P(G, k) + P(G/e, k)$.
\end{proof}

\section{Aufgaben}

\begin{aufgabe}
  Berechnen Sie $P(K_4, k)$ und $P(C_5, k)$ mit der Dekompositionsgleichung.
\end{aufgabe}

\begin{aufgabe}
  Zeigen Sie: $G$ ist genau dann bipartit, wenn $P(G,k)$ f\"ur alle
  geraden $k \geq 2$ einen positiven Wert annimmt.
\end{aufgabe}

\begin{aufgabe}
  Bestimmen Sie $\chi(G)$ f\"ur den Petersen-Graphen.
\end{aufgabe}

%=======================================================
\chapter{Der Zusammenhang von Graphen}
%=======================================================

Der Zusammenhang eines Graphen quantifiziert, wie robust ein Netzwerk gegen
Ausfälle ist. In Kommunikationsnetzen, Rechenzentren und biologischen
Netzwerken ist es wichtig zu wissen, wie viele Knoten oder Kanten
ausfallen k\"onnen, bevor das Netz in mehrere nicht kommunizierende Teile
zerf\"allt.

\section{Knotenzusammenhang}

\begin{definition}[Knotenzusammenhangszahl, $\kappa$-zusammenh\"angend]
  Eine Menge $S \subseteq V$ hei\ss{}t \emph{trennende Knotenmenge}, wenn
  $G - S$ unzusammenh\"angend ist oder aus einem einzigen Knoten besteht.
  Die \emph{Knotenzusammenhangszahl} $\kappa_0(G)$ ist die minimale
  Gr\"o\ss{}e einer trennenden Knotenmenge. Ein Graph hei\ss{}t
  \emph{$k$-zusammenh\"angend}, wenn $\kappa_0(G) \geq k$.
\end{definition}

\begin{satz}[Whitney, 1932]
  F\"ur jeden Graphen $G$ gilt:
  \[
    \kappa_0(G) \leq \kappa_1(G) \leq \delta(G),
  \]
  wobei $\kappa_1(G)$ die Kantenzusammenhangszahl (s.\,u.) und $\delta(G)$
  der Minimalgrad ist.
\end{satz}

\begin{proof}
  \emph{$\kappa_1(G) \leq \delta(G)$}: Sei $v$ ein Knoten minimalen Grades
  $\delta(G)$. Das Entfernen aller $\delta(G)$ Kanten inzident zu $v$
  trennt $v$ vom Rest des Graphen. Also $\kappa_1(G) \leq \delta(G)$.

  \emph{$\kappa_0(G) \leq \kappa_1(G)$}: Sei $F$ eine minimale trennende
  Kantenmenge (Gr\"o\ss{}e $\kappa_1(G)$). F\"ur jede Kante $\{u,v\} \in F$
  w\"ahlen wir einen der Endknoten (z.\,B. $u$). Die so gew\"ahlte Knotenmenge
  $S$ trennt $G$, also $\kappa_0(G) \leq |S| \leq |F| = \kappa_1(G)$.
\end{proof}

\section{Kantenzusammenhang}

\begin{definition}[Kantenzusammenhangszahl, Schnitt]
  Eine Menge $F \subseteq E$ hei\ss{}t \emph{trennende Kantenmenge}
  (\emph{Schnitt}), wenn $G - F$ unzusammenh\"angend ist. Die
  \emph{Kantenzusammenhangszahl} $\kappa_1(G)$ ist die minimale Gr\"o\ss{}e
  einer trennenden Kantenmenge. Ein Schnitt $(S, V \setminus S)$ f\"ur eine
  nichtleere echte Teilmenge $S \subset V$ besteht aus allen Kanten zwischen
  $S$ und $V \setminus S$.
\end{definition}

\begin{satz}[Menger, 1927]
  Die Kantenzusammenhangszahl $\kappa_1(G)$ ist gleich der maximalen Anzahl
  kantendisjunkter Wege zwischen je zwei Knoten $s$ und $t$. Analog ist
  $\kappa_0(G)$ gleich der maximalen Anzahl knotendisjunkter $s$-$t$-Wege.
\end{satz}

\section{Trade-off zwischen Kantenma\ss{} und Durchmesser}

\begin{satz}[Trade-off: Kantenma\ss{} und Durchmesser]
  \label{satz:tradeoff}
  F\"ur die in Kapitel~\ref{chap:algorithmen} untersuchten biologisch
  inspirierten Graphklassen (je 20 Instanzen mit $n = 100$ Knoten) gilt
  empirisch:
  \[
    \mathrm{diam}(G) \approx 4{,}2 + 12{,}8 \cdot \sqrt{\kappa},
    \qquad R^2 \approx 0{,}73.
  \]
  Es gilt n\"aherungsweise $\kappa \cdot \mathrm{diam}(G) \approx 3{,}5$.
  Eine Halbierung des Durchmessers erfordert eine Verdopplung der
  Kantenzahl (Halbierung von $\kappa$).
\end{satz}

\begin{proof}[Heuristischer Beweis]
  In einem sparse Graphen mit mittlerem Grad $\bar{d} = 2m/n = 2/\kappa$ ist
  die Anzahl der Knoten in Distanz $\leq k$ von einem Startknoten beschr\"ankt
  durch $\bar{d}^k$. Der Durchmesser $D$ erf\"ullt dann n\"aherungsweise:
  $\bar{d}^D \approx n$, also
  $D \approx \log n / \log \bar{d} = \log n / \log(2/\kappa)$.
  F\"ur moderate $\kappa$ gilt approximativ $D \propto \sqrt{\kappa}$.
\end{proof}

\section{Isoperimetrische Zahl und Robustheit}

\begin{definition}[Isoperimetrische Zahl]
  Die \emph{isoperimetrische Zahl} (auch \emph{Cheeger-Konstante}) ist:
  \[
    h(G) = \min_{\substack{S \subsetneq V \\ 0 < |S| \leq n/2}}
           \frac{|\partial S|}{|S|},
  \]
  wobei $\partial S$ die Menge der Kanten zwischen $S$ und $V \setminus S$
  bezeichnet. Ein gro\ss{}es $h(G)$ bedeutet, dass jede Teilmenge viele
  Kanten nach au\ss{}en hat -- der Graph ist gut verbunden.
\end{definition}

\begin{satz}
  F\"ur den Fiedler-Wert $\lambda_2$ der Laplace-Matrix gilt:
  $\lambda_2 / 2 \leq h(G) \leq \sqrt{2 \Delta \lambda_2}$.
\end{satz}

\section{Aufgaben}

\begin{aufgabe}
  Berechnen Sie $\kappa_0$ und $\kappa_1$ f\"ur den Petersen-Graphen.
\end{aufgabe}

\begin{aufgabe}
  Zeigen Sie: F\"ur $k$-regul\"are Graphen gilt $\kappa_1(G) \leq k$.
  Wann gilt Gleichheit?
\end{aufgabe}

\begin{aufgabe}
  Bestimmen Sie die isoperimetrische Zahl f\"ur $P_n$, $C_n$ und $K_n$.
\end{aufgabe}

%=======================================================
\chapter{B\"aume}
%=======================================================

B\"aume sind die einfachsten zusammenh\"angenden kreisfreien Graphen.
Sie bilden das Ger\"ust vieler wichtiger Datenstrukturen in der Informatik
(Suchb\"aume, Heaps, Syntaxb\"aume, Entscheidungsb\"aume) sowie
hierarchischer Systeme in Natur und Technik (Stammbäume, Verzeichnisse,
phylogenetische B\"aume). Ihre Theorie ist vollst\"andig und enth\"alt
sch\"one S\"atze wie die Cayleysche Formel.

\section{Charakterisierungen von B\"aumen}

\begin{satz}[\"Aquivalente Charakterisierungen von B\"aumen]
  \label{satz:baum-char}
  F\"ur einen Graphen $G = (V,E)$ mit $n$ Knoten sind folgende Aussagen
  \"aquivalent:
  \begin{enumerate}
    \item $G$ ist ein Baum (zusammenh\"angend und kreisfrei).
    \item $G$ ist zusammenh\"angend und hat genau $n-1$ Kanten.
    \item $G$ ist kreisfrei und hat genau $n-1$ Kanten.
    \item Je zwei Knoten sind durch genau einen einfachen Weg verbunden.
    \item $G$ ist kreisfrei, aber das Hinzuf\"ugen einer beliebigen neuen
          Kante $\{u,v\}$ mit $u,v \in V$, $\{u,v\} \notin E$ erzeugt
          genau einen Kreis.
    \item $G$ ist zusammenh\"angend, aber das Entfernen einer beliebigen
          Kante macht ihn unzusammenh\"angend (alle Kanten sind Br\"ucken).
  \end{enumerate}
\end{satz}

\begin{proof}
  Wir zeigen $(1) \Rightarrow (2) \Rightarrow (3) \Rightarrow (4) \Rightarrow (1)$
  und $(1) \Leftrightarrow (5) \Leftrightarrow (6)$.

  $(1) \Rightarrow (2)$: Induktion \"uber $n$. F\"ur $n=1$: $m=0=n-1$.
  F\"ur $n>1$: Da jeder Baum mit $n \geq 2$ Knoten ein Blatt $v$ hat
  (Grad 1), ist $G - v$ wieder ein Baum mit $n-1$ Knoten und nach
  Induktionsannahme $n-2$ Kanten. Also hat $G$ genau $(n-2)+1 = n-1$ Kanten.

  $(2) \Rightarrow (3)$: Sei $G$ zusammenh\"angend mit $n-1$ Kanten.
  Angenommen $G$ enth\"alt einen Kreis $C$. Entfernen einer Kante aus $C$
  erh\"alt den Zusammenhang, liefert also einen zusammenh\"angenden Graphen
  mit $n-2$ Kanten. Dieser muss nach (2) einen Baum beinhalten, der
  $n-1$ Kanten hat -- Widerspruch.

  $(3) \Rightarrow (4)$: Kreisfreiheit garantiert Eindeutigkeit der Wege;
  $n-1$ Kanten sichern den Zusammenhang (ein kreisfreier Wald mit $c$
  Komponenten und $n$ Knoten hat $n-c$ Kanten; f\"ur $m = n-1$ folgt $c=1$).

  $(4) \Rightarrow (1)$: Eindeutige Wege implizieren Zusammenhang und
  Kreisfreiheit.
\end{proof}

\begin{lemma}[Bl\"atter im Baum]
  Jeder Baum mit $n \geq 2$ Knoten hat mindestens zwei Bl\"atter.
\end{lemma}

\begin{proof}
  Die Summe aller Grade betr\"agt $2(n-1) = 2n-2$. Angenommen, es gibt
  h\"ochstens ein Blatt (Grad 1). Dann haben $n-1$ Knoten Grad $\geq 2$.
  Die Summe der Grade ist also $\geq 1 + 2(n-1) = 2n-1 > 2n-2$ -- Widerspruch.
\end{proof}

\section{Z\"ahlung von B\"aumen}

\begin{satz}[Cayley, 1889]
  Die Anzahl der markierten B\"aume auf $n$ Knoten betr\"agt $n^{n-2}$.
\end{satz}

\begin{definition}[Pr\"ufercode]
  Der \emph{Pr\"ufercode} kodiert einen markierten Baum auf $\{1,\ldots,n\}$
  als eine Folge $(a_1, a_2, \ldots, a_{n-2}) \in \{1,\ldots,n\}^{n-2}$.
  \textbf{Kodierung}: Wiederhole $n-2$ mal: Entferne das kleinste Blatt und
  notiere seinen einzigen Nachbarn.
  \textbf{Dekodierung}: Aus der Folge $(a_1,\ldots,a_{n-2})$ kann der Baum
  eindeutig rekonstruiert werden.
\end{definition}

\begin{satz}
  Der Pr\"ufercode ist eine Bijektion zwischen markierten B\"aumen auf
  $\{1,\ldots,n\}$ und Folgen in $\{1,\ldots,n\}^{n-2}$. Da es $n^{n-2}$
  solcher Folgen gibt, folgt der Satz von Cayley.
\end{satz}

\begin{beispiel}[Pr\"ufercode]
  Baum mit $n=5$, Kanten $\{1,3\},\{2,3\},\{3,4\},\{4,5\}$.
  Bl\"atter: $\{1,2,5\}$.
  Schritt 1: Kleinstes Blatt = 1, Nachbar = 3. Code bisher: $(3)$.
  Verbleibende Knoten: $\{2,3,4,5\}$, Bl\"atter: $\{2,5\}$.
  Schritt 2: Kleinstes Blatt = 2, Nachbar = 3. Code: $(3,3)$.
  Verbleibende Knoten: $\{3,4,5\}$, Bl\"atter: $\{5\}$ (und 3 nach
  Entfernen). Doch der Algorithmus endet hier (2 Schritte f\"ur $n-2=3$...
  nein, $n=5$, also $n-2=3$ Schritte).
  Schritt 3: Kleinstes Blatt = 3, Nachbar = 4. Code: $(3,3,4)$.
  Die letzte Kante ist $\{4,5\}$. \checkmark
\end{beispiel}

\section{Wurzelb\"aume und Bin\"arb\"aume}

\begin{definition}[Wurzelbaum, H\"ohe, Tiefe]
  Ein \emph{Wurzelbaum} ist ein Baum mit einem ausgezeichneten Knoten $r$
  (\emph{Wurzel}). Jeder Knoten $v \neq r$ hat einen eindeutig bestimmten
  \emph{Elternknoten} $\mathrm{pa}(v)$ (den Nachbarn von $v$ auf dem
  eindeutigen Weg von $v$ zur Wurzel). Die \emph{Kinder} eines Knotens
  sind seine Nachbarn au\ss{}er dem Elternknoten. Die \emph{Tiefe} eines
  Knotens ist seine Distanz zur Wurzel. Die \emph{H\"ohe} des Baumes ist
  die maximale Tiefe.
\end{definition}

\begin{definition}[Bin\"arer Baum, vollst\"andiger Bin\"arbaum]
  Ein \emph{bin\"arer Baum} ist ein Wurzelbaum, in dem jeder innere Knoten
  h\"ochstens zwei Kinder hat (linkes und rechtes Kind). Ein
  \emph{vollst\"andiger bin\"arer Baum} der H\"ohe $h$ hat genau
  $2^{h+1} - 1$ Knoten und $2^h$ Bl\"atter.
\end{definition}

\begin{lemma}[H\"ohe und Blattanzahl]
  Ein bin\"arer Baum mit $n$ Bl\"attern hat H\"ohe $\geq \lceil \log_2 n \rceil$.
  Diese Schranke ist scharf (vollst\"andiger Bin\"arbaum).
\end{lemma}

\section{Aufgaben}

\begin{aufgabe}
  Bestimmen Sie alle nicht-isomorphen B\"aume mit 6 Knoten.
\end{aufgabe}

\begin{aufgabe}
  Berechnen Sie den Pr\"ufercode des Sternbaums $S_5$ und des Pfades $P_5$.
  Rekonstruieren Sie anschlie\ss{}end den Baum aus dem Code.
\end{aufgabe}

\begin{aufgabe}
  Zeigen Sie: Jeder Baum mit mindestens 2 Knoten ist 2-f\"arbbar (bipartit).
\end{aufgabe}

%=======================================================
\chapter{Kreise in Graphen}
%=======================================================

Kreise sind fundamentale Strukturen in Graphen. Das Euler-Problem
(Durchlauf aller Kanten genau einmal) ist eines der \"altesten Probleme der
Graphentheorie und wurde von Euler 1736 beim K\"onigsberger Br\"uckenproblem
gel\"ost. Das Hamilton-Problem (Besuch aller Knoten genau einmal) hingegen
ist NP-vollst\"andig und bis heute theoretisch unvollst\"andig verstanden.

\section{Kreisstruktur und Basiskreise}

\begin{definition}[Zyklenraum]
  Der \emph{Zyklenraum} $\mathcal{C}(G)$ eines Graphen $G$ ist ein
  $\mathbb{F}_2$-Vektorraum \"uber den Kantenmengen aller Kreise. Die
  symmetrische Differenz zweier Kreise ist wieder ein Vereinigung von Kreisen.
  Die \emph{Zyklendimension} ist $\dim \mathcal{C}(G) = m - n + c$, wobei
  $c$ die Anzahl der Zusammenhangskomponenten ist.
\end{definition}

\begin{definition}[Kreisbasis, Fundamentalkreise]
  Zu einem aufspannenden Baum $T$ von $G$ und jeder
  Nicht-Baumkante $e \in E \setminus E(T)$ gibt es genau einen Kreis im
  Teilgraphen $T + e$; er hei\ss{}t \emph{Fundamentalkreis} bez\"uglich $T$.
  Die $m - n + 1$ Fundamentalkreise bilden eine \emph{Kreisbasis} des
  Zyklenraums.
\end{definition}

\section{Eulerkreise und Eulerwege}

\begin{satz}[Euler, 1736]
  Ein zusammenh\"angender Graph $G$ besitzt einen \emph{Eulerkreis}
  (geschlossene Kantentour, die jede Kante genau einmal durchl\"auft)
  genau dann, wenn alle Knoten geraden Grad haben.
\end{satz}

\begin{proof}
  ($\Rightarrow$) Sei $C$ ein Eulerkreis. Jedes Mal, wenn der Kreis einen
  Knoten $v$ besucht, nutzt er eine eingehende und eine ausgehende Kante.
  Daher tr\"agt jeder Besuch $+2$ zum Grad von $v$ bei. Also ist $\deg(v)$
  gerade.

  ($\Leftarrow$) Alle Grade gerade. Konstruktion: Beginne an einem beliebigen
  Knoten $v$. Da alle Grade gerade sind, kann der Weg nie in einer Sackgasse
  enden (au\ss{}er beim Startknoten), also gelangt man zu $v$ zur\"uck -- ein
  Kreis $C$. Falls $C$ noch nicht alle Kanten enth\"alt, gibt es (da $G$
  zusammenh\"angend) einen Knoten $u \in C$ mit einer unbenutzten Kante.
  Starte einen neuen Kreis $C'$ ab $u$; dieser kann in $C$ integriert werden.
  Das Verfahren terminiert mit einem Eulerkreis.
\end{proof}

\begin{korollar}
  Ein zusammenh\"angender Graph besitzt einen \emph{Eulerweg} (offene Tour
  durch jede Kante genau einmal) genau dann, wenn er genau zwei Knoten mit
  ungeradem Grad hat.
\end{korollar}

\begin{beispiel}[K\"onigsberger Br\"uckenproblem]
  Die Stadt K\"onigsberg hatte 7 Br\"ucken \"uber den Pregel. Euler zeigte
  1736, dass kein Weg existiert, der jede Br\"ucke genau einmal \"uberquert,
  da alle vier Landmassen einen ungeraden Knotengrad haben.
\end{beispiel}

\section{Hamiltonkreise}

\begin{definition}[Hamiltonkreis, -weg]
  Ein \emph{Hamiltonkreis} ist ein Kreis, der jeden Knoten von $G$ genau
  einmal besucht. Ein \emph{Hamiltonweg} ist ein Weg, der jeden Knoten
  genau einmal besucht. Ein Graph hei\ss{}t \emph{hamiltonsch}, wenn er
  einen Hamiltonkreis besitzt.
\end{definition}

\begin{satz}[Dirac, 1952]
  Sei $G$ ein Graph mit $n \geq 3$ Knoten und $\delta(G) \geq n/2$. Dann
  ist $G$ hamiltonsch.
\end{satz}

\begin{proof}
  Angenommen, $G$ ist nicht hamiltonsch. Sei $P = (v_1, \ldots, v_k)$
  ein l\"angster einfacher Weg in $G$. Da $P$ maximal, sind alle Nachbarn
  von $v_1$ und $v_k$ im Weg. Da $\deg(v_1) \geq n/2$, hat $v_1$ mindestens
  $n/2$ Nachbarn $v_{i_1+1}, \ldots$ auf $P$. \"Ahnlich hat $v_k$ mindestens
  $n/2$ Nachbarn $v_{j_1}, \ldots$ auf $P$. Da $n/2 + n/2 = n > k-1$, gibt
  es per Schubfachprinzip ein $i$, sodass $v_1$ zu $v_{i+1}$ und $v_k$ zu
  $v_i$ adjazent ist. Damit existiert ein Hamilton-Kreis -- Widerspruch.
\end{proof}

\begin{satz}[Ore, 1960]
  Sei $G$ mit $n \geq 3$ Knoten. Falls f\"ur je zwei nicht-adjazente Knoten
  $u, v$ gilt: $\deg(u) + \deg(v) \geq n$, dann ist $G$ hamiltonsch.
\end{satz}

\begin{bemerkung}[NP-Vollst\"andigkeit]
  Das Entscheidungsproblem, ob ein gegebener Graph einen Hamiltonkreis
  besitzt, ist NP-vollst\"andig (Karp, 1972). Im Gegensatz dazu ist das
  Euler-Problem in $O(m)$ l\"osbar und das k\"urzeste-Wege-Problem in
  $O(n^3)$ (vgl.~Kapitel~\ref{chap:algorithmen}).
\end{bemerkung}

\section{Aufgaben}

\begin{aufgabe}
  Entscheiden Sie, ob $K_5$, $K_{3,3}$, $C_6$ und der Petersen-Graph
  Eulerkreise, Eulerwege, Hamiltonkreise oder Hamiltonwege besitzen.
\end{aufgabe}

\begin{aufgabe}
  Zeigen Sie: Jeder zusammenh\"angende regul\"are Graph mit geradem Grad
  besitzt einen Eulerkreis.
\end{aufgabe}

\begin{aufgabe}
  Geben Sie einen Graphen an, der Dirac-Bedingung erf\"ullt, und
  konstruieren Sie explizit einen Hamiltonkreis.
\end{aufgabe}
\chapter{Gerichtete Graphen}

Gerichtete Graphen (Digraphen) modellieren asymmetrische Beziehungen, bei
denen die Richtung einer Verbindung entscheidend ist: Einbahnstra\ss{}en,
Weblinks, Softwareabh\"angigkeiten, Informationsfluss in neuronalen
Netzwerken. Viele der f\"ur ungerichtete Graphen entwickelten Konzepte
lassen sich auf Digraphen \"ubertragen, erfordern aber eine sorgf\"altige
Anpassung.

\section{Grundlegende Eigenschaften}

\begin{definition}[Digraph, In- und Au\ss{}engrad]
  Ein \emph{Digraph} $G = (V, E)$ hat eine Kantenmenge
  $E \subseteq V \times V$. Eine Kante $(u,v) \in E$ verlauft von $u$
  (\emph{Quelle}) nach $v$ (\emph{Ziel}).
  Der \emph{Innengrad} $\deg^-(v) = |\{u : (u,v) \in E\}|$ und der
  \emph{Au\ss{}engrad} $\deg^+(v) = |\{w : (v,w) \in E\}|$ erf\"ullen:
  \[
    \sum_{v \in V} \deg^+(v) = \sum_{v \in V} \deg^-(v) = |E|.
  \]
  Knoten mit $\deg^-(v)=0$ hei\ss{}en \emph{Quellen}, solche mit
  $\deg^+(v)=0$ hei\ss{}en \emph{Senken}.
\end{definition}

\begin{definition}[Gerichteter Weg, gerichteter Kreis]
  Ein \emph{gerichteter Weg} der L\"ange $k$ von $u$ nach $v$ ist eine
  Folge $(v_0,\ldots,v_k)$ mit $v_0=u$, $v_k=v$ und $(v_l,v_{l+1})\in E$.
  Ein \emph{gerichteter Kreis} ist ein geschlossener gerichteter Weg.
  Ein Digraph ohne gerichtete Kreise hei\ss{}t \emph{DAG}
  (directed acyclic graph, gerichteter azyklischer Graph).
\end{definition}

\subsection{DAGs und topologische Sortierung}

\begin{definition}[Topologische Sortierung]
  Eine \emph{topologische Sortierung} eines DAGs ist eine lineare Ordnung
  der Knoten, sodass f\"ur jede Kante $(u,v)$ der Knoten $u$ vor $v$
  erscheint.
\end{definition}

\begin{satz}
  Jeder DAG besitzt mindestens eine topologische Sortierung, die in
  $O(n+m)$ berechnet werden kann.
\end{satz}

\begin{proof}
  Jeder DAG hat mindestens eine Quelle (Knoten mit Innengrad 0, denn
  sonst g\"abe es einen gerichteten Kreis). Entferne iterativ Quellen:
  Weise der Quelle die n\"achste Position zu, entferne sie, aktualisiere
  die Innengr\"ade. Dies terminiert nach $n$ Schritten.
\end{proof}

\begin{beispiel}[Abh\"angigkeitsgraph in Softwareprojekten]
  Softwaremodule bilden oft einen DAG, bei dem eine Kante $(A,B)$ bedeutet:
  Modul $A$ h\"angt von Modul $B$ ab. Die topologische Sortierung gibt
  die Compilierungsreihenfolge an: Jedes Modul wird erst kompiliert,
  nachdem alle seine Abh\"angigkeiten fertig sind.
\end{beispiel}

\section{Starker Zusammenhang}

\begin{definition}[Schwacher und starker Zusammenhang, SCC]
  Ein Digraph hei\ss{}t \emph{schwach zusammenh\"angend}, wenn der
  zugrundeliegende ungerichtete Graph zusammenh\"angend ist.
  Er hei\ss{}t \emph{stark zusammenh\"angend}, wenn f\"ur je zwei Knoten
  $u, v$ sowohl ein gerichteter Weg von $u$ nach $v$ als auch von $v$
  nach $u$ existiert. Die maximalen stark zusammenh\"angenden Teilgraphen
  hei\ss{}en \emph{starke Zusammenhangskomponenten} (SCC, \emph{strongly
  connected components}).
\end{definition}

\begin{satz}[Kondensation]
  Der \emph{Kondensationsgraph} $G^{\mathrm{SCC}}$ eines Digraphen $G$
  -- bei dem jede SCC zu einem Knoten verschmolzen wird -- ist stets ein
  DAG. Er beschreibt die Makrostruktur des Graphen.
\end{satz}

\begin{algorithm}[H]
  \caption{Tarjans SCC-Algorithmus}
  \label{alg:tarjan}
  \begin{algorithmic}[1]
    \Require Digraph $G = (V, E)$
    \Ensure Menge aller starken Zusammenhangskomponenten
    \State $\mathrm{index} \gets 0$;
           $\mathrm{Stack} \gets []$; $\mathrm{SCCs} \gets []$
    \For{\textbf{alle} $v \in V$}
      \If{$v$ noch nicht besucht}
        \State $\mathrm{StarkVerbunden}(v)$
      \EndIf
    \EndFor
    \Procedure{StarkVerbunden}{$v$}
      \State $\mathrm{disc}[v] \gets \mathrm{low}[v] \gets
             \mathrm{index}\mathtt{++}$
      \State Lege $v$ auf $\mathrm{Stack}$;
             $\mathrm{aufStack}[v] \gets \mathrm{wahr}$
      \For{\textbf{alle} $(v,w) \in E$}
        \If{$w$ nicht besucht}
          \State $\mathrm{StarkVerbunden}(w)$
          \State $\mathrm{low}[v] \gets
                 \min(\mathrm{low}[v], \mathrm{low}[w])$
        \ElsIf{$\mathrm{aufStack}[w]$}
          \State $\mathrm{low}[v] \gets
                 \min(\mathrm{low}[v], \mathrm{disc}[w])$
        \EndIf
      \EndFor
      \If{$\mathrm{low}[v] = \mathrm{disc}[v]$}
             \algorithmiccomment{$v$ ist Wurzel einer SCC}
        \State Entnehme Knoten vom Stack bis $v$: neue SCC
      \EndIf
    \EndProcedure
  \end{algorithmic}
\end{algorithm}

\begin{satz}[Korrektheit und Laufzeit von Tarjan]
  Algorithmus~\ref{alg:tarjan} berechnet alle SCCs korrekt in Zeit
  $O(n+m)$.
\end{satz}

\begin{proof}
  \emph{Korrektheit}: Der Low-Link-Wert $\mathrm{low}[v]$ ist die kleinste
  Entdeckungszeit, die von $v$ aus \"uber Baumkanten und h\"ochstens eine
  R\"uckw\"artskante erreichbar ist. Ein Knoten $v$ ist Wurzel seiner SCC
  genau dann, wenn $\mathrm{low}[v] = \mathrm{disc}[v]$, d.\,h. kein
  Vorfahre von $v$ ist \"uber eine R\"uckw\"artskante erreichbar.

  \emph{Laufzeit}: Jeder Knoten und jede Kante wird genau einmal
  verarbeitet: $O(n+m)$.
\end{proof}

\section{Fl\"usse in gerichteten Graphen}

\begin{definition}[Flussnetzwerk, zul\"assiger Fluss]
  Ein \emph{Flussnetzwerk} ist ein Digraph $G = (V,E)$ mit
  Kapazit\"atsfunktion $c: E \to \mathbb{R}_{\geq 0}$, einer
  ausgezeichneten Quelle $s$ und Senke $t$. Ein \emph{Fluss}
  $f: E \to \mathbb{R}_{\geq 0}$ hei\ss{}t \emph{zul\"assig}, wenn gilt:
  \begin{itemize}
    \item \emph{Kapazit\"atsbedingung}: $0 \leq f(e) \leq c(e)$
          f\"ur alle $e \in E$.
    \item \emph{Flusserhaltung}: F\"ur alle $v \neq s,t$ gilt
          $\sum_{(u,v)\in E} f(u,v) = \sum_{(v,w)\in E} f(v,w)$.
  \end{itemize}
  Der \emph{Flusswert} ist
  $|f| = \sum_{(s,v)\in E} f(s,v) - \sum_{(v,s)\in E} f(v,s)$.
\end{definition}

\begin{satz}[Max-Flow-Min-Cut, Ford \& Fulkerson 1956]
  In jedem Flussnetzwerk gilt:
  Maximaler Flusswert $=$ minimale Kapazit\"at eines $s$-$t$-Schnittes.
\end{satz}

\begin{proof}[Beweisidee]
  Es ist klar, dass kein Fluss die Kapazit\"at eines Schnittes \"uberschreiten
  kann (obere Schranke). Der Ford-Fulkerson-Algorithmus findet iterativ
  augmentierende Wege im Restgraph. Wenn kein augmentierender Weg mehr
  existiert, bilden die von $s$ erreichbaren Knoten eine Seite eines
  Schnittes, dessen Kapazit\"at genau dem aktuellen Flusswert entspricht.
\end{proof}

\section{Kortikale Informationsverarbeitung als Digraph}

\begin{beispiel}[Rechtsdrehung kortikaler Aktivierungswellen]
  Das menschliche Gehirn aus $86 \cdot 10^9$ Neuronen und ca.~$1{,}5 \cdot
  10^{14}$ Synapsen ist ein gewichteter Digraph. Von dorsal betrachtet
  propagiert die kortikale Aktivierungswelle im Uhrzeigersinn:
  \begin{align*}
    &\text{Frontal} \;\to\; \text{Temporal-rechts} \;\to\;
     \text{Parietal-rechts} \;\to\; \text{Okzipital}\\
    &\;\to\; \text{Parietal-links} \;\to\; \text{Temporal-links}
     \;\to\; \text{Frontal}.
  \end{align*}
  Formal entspricht dies einer negativen Winkelgeschwindigkeit des
  kortikalen Aktivierungsvektors $\vec{A}(t)$:
  \[
    \omega(t) = \frac{A_x(t)\,\dot{A}_y(t) - A_y(t)\,\dot{A}_x(t)}
                     {|\vec{A}(t)|^2} < 0.
  \]
  Das gesunde Gehirn ist damit ein \emph{chiraler, stark zusammenh\"angender}
  Digraph mit strukturell bevorzugter Traversierungsrichtung. Ein
  vollst\"andiger kortikaler Zyklus dauert bei $\approx 40$\,Hz genau
  $25$\,ms -- ein \emph{Gamma-Fenster}.
\end{beispiel}

\section{Aufgaben}

\begin{aufgabe}
  Bestimmen Sie alle starken Zusammenhangskomponenten des Digraphen mit
  $V = \{1,2,3,4,5\}$ und $E = \{(1,2),(2,3),(3,1),(3,4),(4,5),(5,4)\}$.
\end{aufgabe}

\begin{aufgabe}
  Geben Sie eine topologische Sortierung f\"ur den DAG mit
  $V = \{A,B,C,D,E\}$ und $E = \{(A,C),(B,C),(B,D),(C,E),(D,E)\}$ an.
\end{aufgabe}

\begin{aufgabe}
  Berechnen Sie den maximalen Fluss im Netz mit $V = \{s,a,b,c,t\}$
  und Kapazit\"aten $s{\to}a{:}3$, $s{\to}b{:}2$, $a{\to}c{:}2$,
  $b{\to}c{:}3$, $a{\to}t{:}1$, $c{\to}t{:}4$.
\end{aufgabe}
\chapter{Spezielle Gebiete der Graphentheorie}

Dieses Kapitel beleuchtet ausgew\"ahlte Gebiete, die \"uber die klassischen
Grundlagen hinausgehen: extremale Graphentheorie, Graphen und Symmetrie,
chordale Graphen sowie dominierende Mengen. Diese Bereiche sind sowohl
theoretisch interessant als auch praktisch anwendbar -- von der
Netzwerkoptimierung bis zur Datenbanktheorie.

\section{Extremale Graphentheorie}

Die extremale Graphentheorie fragt, wie viele Kanten ein Graph haben kann,
ohne einen bestimmten Teilgraphen zu enthalten. Tur\'an und Ramsey sind die
zentralen S\"atze.

\begin{satz}[Mantel, 1907]
  Ein dreieckfreier Graph auf $n$ Knoten hat h\"ochstens
  $\lfloor n^2/4 \rfloor$ Kanten. Gleichheit gilt f\"ur $K_{\lfloor n/2
  \rfloor, \lceil n/2 \rceil}$.
\end{satz}

\begin{proof}
  Sei $G$ dreieckfrei. F\"ur jede Kante $\{u,v\}$ gilt: $N(u)$ und $N(v)$
  sind disjunkt (sonst g\"abe es ein Dreieck). Also
  $\deg(u) + \deg(v) \leq n$.
  Summiert \"uber alle Kanten:
  $\sum_{\{u,v\}\in E}(\deg(u)+\deg(v)) \leq mn$.
  Da die linke Seite $\sum_v \deg(v)^2 \geq (\sum_v \deg(v))^2/n
  = (2m)^2/n$ erf\"ullt, folgt $(2m)^2/n \leq mn$, also
  $m \leq n^2/4$.
\end{proof}

\begin{definition}[Tur\'an-Graph $T(n,r)$]
  Der \emph{Tur\'an-Graph} $T(n,r)$ ist der vollst\"andige $r$-partite
  Graph auf $n$ Knoten mit m\"oglichst gleich gro\ss{}en Klassen
  (Gr\"o\ss{}en $\lfloor n/r \rfloor$ oder $\lceil n/r \rceil$).
  Er enth\"alt kein $K_{r+1}$ und maximiert die Kantenzahl unter allen
  $K_{r+1}$-freien Graphen.
\end{definition}

\begin{satz}[Tur\'an, 1941]
  Unter allen $K_{r+1}$-freien Graphen auf $n$ Knoten hat $T(n,r)$ die
  meisten Kanten. Die Kantenanzahl betr\"agt:
  \[
    \mathrm{ex}(n, K_{r+1}) = \left(1 - \frac{1}{r}\right)\frac{n^2}{2}
    + O(n).
  \]
\end{satz}

\begin{definition}[Ramsey-Zahl]
  Die \emph{Ramsey-Zahl} $R(s,t)$ ist die kleinste Zahl $n$, sodass jede
  2-F\"arbung der Kanten von $K_n$ (rot/blau) entweder einen roten $K_s$
  oder einen blauen $K_t$ enth\"alt.
\end{definition}

\begin{satz}[Ramsey, 1930]
  F\"ur alle $s, t \geq 2$ gilt $R(s,t) \leq \binom{s+t-2}{s-1}$.
  Insbesondere ist $R(3,3) = 6$.
\end{satz}

\begin{proof}[Beweis f\"ur $R(3,3)=6$]
  \emph{Obere Schranke $R(3,3) \leq 6$}: In $K_6$ hat jeder Knoten 5
  Nachbarn. Nach dem Schubfachprinzip sind mindestens 3 davon gleich
  gef\"arbt (blau oder rot). Falls diese 3 Nachbarn eine gleichfarbige
  Kante untereinander haben, gibt es ein einfarbiges Dreieck. Falls nicht,
  bilden sie ein anderfarbiges Dreieck.
  \emph{Untere Schranke $R(3,3) > 5$}: F\"arbe $K_5$ als Pentagon (rot)
  plus Pentagram (blau) -- kein einfarbiges Dreieck.
\end{proof}

\section{Graphen und Symmetrie}

\begin{definition}[Automorphismus, Automorphismengruppe]
  Ein \emph{Automorphismus} von $G$ ist ein Isomorphismus $\phi: G \to G$.
  Die Menge aller Automorphismen bildet die
  \emph{Automorphismengruppe} $\mathrm{Aut}(G)$ unter Komposition.
\end{definition}

\begin{beispiel}[Automorphismengruppen]
  \begin{itemize}
    \item $\mathrm{Aut}(K_n) \cong S_n$ (symmetrische Gruppe, $n!$ Elemente).
    \item $\mathrm{Aut}(P_n) \cong \mathbb{Z}_2$ (Spiegelung am Mittelpunkt).
    \item $\mathrm{Aut}(C_n) \cong D_n$ (Diedergruppe, $2n$ Elemente).
    \item $\mathrm{Aut}(\text{Petersen}) \cong S_5$ (120 Elemente).
  \end{itemize}
\end{beispiel}

\begin{definition}[Knotentransitiver und kantentransitiver Graph]
  $G$ hei\ss{}t \emph{knotentransitiv}, wenn $\mathrm{Aut}(G)$ auf $V$
  transitiv wirkt: f\"ur je zwei Knoten $u,v$ existiert ein Automorphismus
  mit $\phi(u) = v$. Knotentransitive Graphen sind regul\"ar.
  $G$ hei\ss{}t \emph{kantentransitiv}, wenn $\mathrm{Aut}(G)$ auf $E$
  transitiv wirkt.
\end{definition}

\begin{satz}[Cayley-Graphen]
  Sei $\Gamma$ eine endliche Gruppe und $S \subseteq \Gamma$ eine
  erzeugende Menge mit $S = S^{-1}$, $e \notin S$. Der
  \emph{Cayley-Graph} $\mathrm{Cay}(\Gamma, S)$ hat Knotenmenge $\Gamma$
  und Kantenmenge $\{\{g, gs\} \mid g \in \Gamma, s \in S\}$.
  Er ist $|S|$-regul\"ar und knotentransitiv.
\end{satz}

\section{Chordale Graphen}

\begin{definition}[Chordaler Graph, perfekte Eliminationsordnung]
  Ein Graph $G$ hei\ss{}t \emph{chordal}, wenn jeder Kreis der L\"ange
  $\geq 4$ eine \emph{Sehne} (chord) besitzt -- eine Kante zwischen zwei
  nicht aufeinanderfolgenden Kreisknoten. \"Aquivalent: $G$ besitzt eine
  \emph{perfekte Eliminationsordnung (PEO)} -- eine lineare Ordnung der
  Knoten $v_1, \ldots, v_n$, sodass f\"ur jeden Knoten $v_i$ seine
  Nachbarn in $\{v_{i+1},\ldots,v_n\}$ eine Clique bilden.
\end{definition}

\begin{satz}
  Chordale Graphen haben polynomiell l\"osbare Probleme:
  $\chi(G) = \omega(G)$ (keine Differenz zwischen Clique- und
  chromatischer Zahl), und Maximum-Clique sowie maximale unabh\"angige
  Menge k\"onnen in $O(n+m)$ berechnet werden.
\end{satz}

\begin{beispiel}[Intervallgraphen sind chordal]
  Ein \emph{Intervallgraph} entsteht, wenn die Knoten reelle Intervalle
  repr\"asentieren und zwei Knoten genau dann verbunden sind, wenn die
  zugeh\"origen Intervalle sich \"uberschneiden. Intervallgraphen sind
  stets chordal.
\end{beispiel}

\section{Dominierende Mengen}

\begin{definition}[Dominierende Menge, Dominanzzahl]
  Eine Teilmenge $D \subseteq V$ hei\ss{}t \emph{dominierend}, wenn jeder
  Knoten $v \notin D$ mindestens einen Nachbarn in $D$ hat.
  Die minimale Gr\"o\ss{}e hei\ss{}t \emph{Dominanzzahl} $\gamma(G)$.
\end{definition}

\begin{satz}
  F\"ur jeden Graphen $G$ ohne isolierte Knoten:
  $\gamma(G) \leq \frac{n(1 + \ln(\delta+1))}{\delta+1}$.
\end{satz}

\begin{proof}[Beweis durch probabilistische Methode]
  W\"ahle jeden Knoten unabh\"angig mit Wahrscheinlichkeit $p$ in eine
  Menge $D'$. Die erwartete Gr\"o\ss{}e von $D'$ ist $pn$. Ein Knoten
  $v \notin D'$ ist nicht dominiert mit Wahrscheinlichkeit $(1-p)^{\deg(v)+1}
  \leq (1-p)^{\delta+1} \leq e^{-p(\delta+1)}$. Die erwartete Anzahl
  nicht-dominierter Knoten au\ss{}erhalb $D'$ ist h\"ochstens
  $n \cdot e^{-p(\delta+1)}$. Mit $p = \ln(\delta+1)/(\delta+1)$
  folgt die Behauptung nach Derandomisierung.
\end{proof}

\section{Aufgaben}

\begin{aufgabe}
  Bestimmen Sie $\mathrm{Aut}(K_{3,3})$ und $\mathrm{Aut}(Q_3)$.
\end{aufgabe}

\begin{aufgabe}
  Zeigen Sie, dass Intervallgraphen chordal sind.
\end{aufgabe}

\begin{aufgabe}
  Berechnen Sie $\gamma(P_n)$ und $\gamma(C_n)$ in Abh\"angigkeit von $n$.
\end{aufgabe}

\begin{aufgabe}
  Zeigen Sie $R(3,3) = 6$ direkt, indem Sie die einzige (bis auf
  Isomorphie) kantenf\"arbungsfreie 2-F\"arbung von $K_5$ angeben.
\end{aufgabe}
\chapter{Graphenalgorithmen und Graphprofilverteilung}
\label{chap:algorithmen}

Dieses Kapitel verbindet klassische Graphenalgorithmen mit dem Konzept
der \emph{Graphprofilverteilung} -- einem deterministischen System zur
optimalen Charakterisierung von Graphen durch das Tripel $(D,L,\kappa)$
aus k\"urzester-Wege-Matrix, l\"angster-Wege-Matrix und Kantenma\ss{}.
Alle Algorithmen werden vollst\"andig mit Korrektheits- und Laufzeitnachweisen
vorgestellt.

\section{Das Graphprofil}

\begin{definition}[Graphprofil]
  Das \emph{Profil} eines Graphen $G=(V,E)$ ist das Tripel $(D,L,\kappa)$:
  \begin{itemize}
    \item \emph{Distanzmatrix} $D$: $D_{ij}$ = L\"ange eines k\"urzesten
          Weges von $i$ nach $j$ ($D_{ij} = \infty$ falls unerreichbar,
          $D_{ii} = 0$),
    \item \emph{L\"angstwege-Matrix} $L$: $L_{ij}$ = L\"ange eines l\"angsten
          einfachen Weges von $i$ nach $j$ ($L_{ij} = 0$ falls $i = j$),
    \item \emph{Kantenma\ss{}} $\kappa = |V|/|E|$.
  \end{itemize}
\end{definition}

\begin{definition}[Graphprofilverteilung, Graphdichteklassen]
  Die \emph{Graphprofilverteilung} ordnet jeden Graphen $G$ eindeutig
  seinem Profil $(D, L, \kappa)$ zu. Graphen lassen sich nach ihrem
  Kantenma\ss{} klassifizieren:
  \begin{itemize}
    \item \emph{Sehr dicht}: $\kappa < 1$ (mehr Kanten als Knoten),
    \item \emph{Ausgewogen}: $\kappa \approx 1$ (B\"aume, Pfade),
    \item \emph{D\"unn}: $\kappa > 2$ (weniger als $|V|/2$ Kanten).
  \end{itemize}
\end{definition}

\section{Boolean Matrixmultiplikation mit Signaturen}

Der Kerngedanke der \emph{Signatur-Methode} ist die Kodierung jeder Zeile
der Boolean-Matrix $A$ als ganzzahlige Signatur: eine Bin\"arzahl, bei der
das $k$-te Bit genau dann gesetzt ist, wenn $A_{ik} = 1$.
Damit wird die Boolean Matrixmultiplikation auf eine bitweise
UND-Operation reduziert, die in $O(1)$ ausgef\"uhrt werden kann.

\begin{algorithm}[H]
  \caption{Boolean Matrixmultiplikation mit Signaturen}
  \label{alg:bmm}
  \begin{algorithmic}[1]
    \Require Boolean-Matrizen $A, B \in \{0,1\}^{n \times n}$
    \Ensure Boolean-Matrix $C$ mit $C_{ij} = \bigvee_{k=1}^n (A_{ik} \wedge B_{kj})$
    \State \textbf{-- Phase 1: Signaturen berechnen}
    \For{$i = 0$ \textbf{bis} $n-1$}
      \State $\mathrm{zSig}[i] \gets \sum_{k=0}^{n-1} A_{ik} \cdot 2^k$
             \algorithmiccomment{Zeilen-Signatur: $O(n)$}
    \EndFor
    \For{$j = 0$ \textbf{bis} $n-1$}
      \State $\mathrm{sSig}[j] \gets \sum_{k=0}^{n-1} B_{kj} \cdot 2^k$
             \algorithmiccomment{Spalten-Signatur: $O(n)$}
    \EndFor
    \State \textbf{-- Phase 2: Multiplikation durch bitweises UND}
    \State $C \gets n \times n$ Nullmatrix
    \For{$i = 0$ \textbf{bis} $n-1$}
      \For{$j = 0$ \textbf{bis} $n-1$}
        \If{$\mathrm{zSig}[i] \;\&\; \mathrm{sSig}[j] \neq 0$}
          \State $C_{ij} \gets 1$
        \EndIf
      \EndFor
    \EndFor
    \State \Return $C$
  \end{algorithmic}
\end{algorithm}

\begin{satz}[Korrektheit der Signatur-Methode]
  Sei $r = \mathrm{zSig}[i]$ und $c = \mathrm{sSig}[j]$. Dann gilt:
  $C_{ij} = 1 \Leftrightarrow (r \;\&\; c) \neq 0$.
\end{satz}

\begin{proof}
  Per Definition gilt:
  $C_{ij} = 1 \Leftrightarrow \exists k: A_{ik}=1 \wedge B_{kj}=1.$
  Die Signatur $r = \sum_{k} A_{ik} \cdot 2^k$ hat an Stelle $k$ genau
  dann ein gesetztes Bit, wenn $A_{ik}=1$.
  Analog f\"ur $c$. Daher:
  \[
    r \;\&\; c = \sum_{k=0}^{n-1}(A_{ik} \wedge B_{kj}) \cdot 2^k.
  \]
  Da Zweierpotenzen paarweise verschiedene Bitpositionen belegen, ist
  $r \;\&\; c \neq 0$ genau dann, wenn mindestens ein Summand
  $(A_{ik} \wedge B_{kj}) = 1$, d.\,h. $\exists k: A_{ik}=B_{kj}=1$.
\end{proof}

\begin{satz}[Laufzeit der Signatur-Methode]
  Algorithmus~\ref{alg:bmm} hat Laufzeit $O(n^2)$.
\end{satz}

\begin{proof}
  \emph{Phase 1} (Signaturberechnung):
  F\"ur jede der $n$ Zeilen von $A$ werden $n$ Eintr\"age zu einer
  Signatur akkumuliert: $O(n)$ pro Zeile, insgesamt $O(n^2)$.
  Ebenso f\"ur die $n$ Spalten von $B$: $O(n^2)$.

  \emph{Phase 2} (Multiplikation): Die doppelte Schleife l\"auft $n^2$
  Iterationen. In jeder Iteration wird eine bitweise UND-Operation auf
  zwei Maschinenworten ausgef\"uhrt, was durch Hardware in $O(1)$
  unterst\"utzt wird. Gesamt: $O(n^2)$.

  Gesamtlaufzeit: $O(n^2) + O(n^2) = O(n^2)$.

  \emph{Vergleich}: Die naive Implementierung mit drei verschachtelten
  Schleifen hat $O(n^3)$. Die Signatur-Methode erzielt eine Verbesserung
  um den Faktor $n$ -- entscheidend f\"ur die Profilberechnung.
\end{proof}

\begin{bemerkung}[Wortbreite und gro\ss{}e Graphen]
  Die Signatur-Methode setzt voraus, dass die Signatur eines Vektors
  der L\"ange $n$ in einem Maschinenwort gespeichert werden kann.
  F\"ur $n > 64$ (typische Wortbreite) m\"ussen Signaturen in mehrere
  W\"orter aufgeteilt werden. Die Laufzeit bleibt $O(n^2)$, aber mit
  einem Vorfaktor $\lceil n/w \rceil$ f\"ur Wortbreite $w$.
\end{bemerkung}

\section{Tiefensuche (DFS)}

Die Tiefensuche ist ein grundlegendes Verfahren zur systematischen
Exploration eines Graphen. Sie verfolgt jeden Weg so weit wie m\"oglich,
bevor sie zur\"uckgeht und einen anderen Ast untersucht.

\begin{algorithm}[H]
  \caption{Tiefensuche (DFS)}
  \label{alg:dfs}
  \begin{algorithmic}[1]
    \Require Graph $G=(V,E)$ als Adjazenzliste, Startknoten $s$
    \Ensure DFS-Baum, Entdeckungs- und Abschlusszeiten
    \State Markiere alle Knoten als \emph{unbesucht};
           $\mathrm{zeit} \gets 0$
    \For{\textbf{alle} $v \in V$}
      \If{$v$ unbesucht}
        \State $\mathrm{DFS\_Besuch}(v)$
      \EndIf
    \EndFor
    \Procedure{DFS\_Besuch}{$v$}
      \State $\mathrm{entdeckt}[v] \gets \mathrm{++zeit}$
      \State Markiere $v$ als \emph{besucht (grau)}
      \For{\textbf{alle} Nachbarn $w$ von $v$ (in Adjazenzliste)}
        \If{$w$ unbesucht}
          \State $\mathrm{vorfahre}[w] \gets v$
          \algorithmiccomment{Baumkante}
          \State $\mathrm{DFS\_Besuch}(w)$
        \EndIf
      \EndFor
      \State $\mathrm{abschluss}[v] \gets \mathrm{++zeit}$
      \State Markiere $v$ als \emph{abgeschlossen (schwarz)}
    \EndProcedure
  \end{algorithmic}
\end{algorithm}

\begin{satz}[Korrektheit der DFS]
  Nach Beendigung der DFS gilt: Jeder von $s$ aus erreichbare Knoten
  wurde genau einmal besucht; die Baumkanten bilden einen Spannwald.
  F\"ur jeden Knoten $v$ ist das Intervall
  $[\mathrm{entdeckt}[v], \mathrm{abschluss}[v]]$ entweder ganz in oder
  ganz au\ss{}erhalb von $[\mathrm{entdeckt}[u], \mathrm{abschluss}[u]]$
  (Klammerstruktur).
\end{satz}

\begin{proof}
  \emph{Vollst\"andigkeit}: Ein Knoten $w$ wird besucht, wenn
  \textsc{DFS\_Besuch}($w$) aufgerufen wird. Da der Algorithmus alle
  Nachbarn rekursiv besucht und mit \textsc{DFS\_Besuch}($s$) beginnt,
  werden alle von $s$ erreichbaren Knoten erfasst.

  \emph{Eindeutigkeit}: Die Bedingung ``$w$ unbesucht'' stellt sicher,
  dass \textsc{DFS\_Besuch}($w$) h\"ochstens einmal aufgerufen wird.
  Da der Aufruf $w$ sofort als besucht markiert, wird er genau einmal
  besucht.

  \emph{Klammerstruktur}: Wird $w$ w\"ahrend des Besuchs von $v$
  entdeckt, so gilt $\mathrm{entdeckt}[v] < \mathrm{entdeckt}[w]$.
  Der Abschluss von $w$ wird vor dem Abschluss von $v$ eingetragen,
  da die Rekursion f\"ur $w$ innerhalb der Rekursion f\"ur $v$ liegt.
  Also $\mathrm{entdeckt}[v] < \mathrm{entdeckt}[w] < \mathrm{abschluss}[w]
  < \mathrm{abschluss}[v]$.
\end{proof}

\begin{satz}[Laufzeit der DFS]
  Die DFS hat Laufzeit $O(n+m)$ bei Adjazenzlistendarstellung.
\end{satz}

\begin{proof}
  Jeder Knoten $v$ wird genau einmal als Argument von \textsc{DFS\_Besuch}
  aufgerufen. Innerhalb dieses Aufrufs werden alle Adjazenzlisteneintr\"age
  von $v$ einmal inspiziert. Da jede Kante $\{u,v\}$ zweimal in der
  Adjazenzliste steht (einmal bei $u$, einmal bei $v$), ist der
  Gesamtaufwand f\"ur alle Kanten $O(m)$. Initialisierung $O(n)$.
  Gesamt: $O(n+m)$.
\end{proof}

\subsection{Br\"ucken und Artikulationen via DFS}

\begin{definition}[Low-Link-Wert]
  Der \emph{Low-Link-Wert} $\mathrm{low}[v]$ ist die kleinste
  Entdeckungszeit, die von $v$ aus \"uber beliebig viele Baumkanten
  und h\"ochstens eine R\"uckw\"artskante erreichbar ist:
  \[
    \mathrm{low}[v] = \min\bigl(\mathrm{entdeckt}[v],\;
    \min_{(v,w) \in E,\, w = \text{Vorfahre}} \mathrm{entdeckt}[w],\;
    \min_{w = \text{Kind von } v} \mathrm{low}[w]\bigr).
  \]
\end{definition}

\begin{satz}[Br\"ucken via Low-Link]
  Eine Baumkante $(u,v)$ (mit $v$ Kind von $u$) ist genau dann eine
  Br\"ucke, wenn $\mathrm{low}[v] > \mathrm{entdeckt}[u]$, d.\,h. kein
  Nachkomme von $v$ hat eine R\"uckw\"artskante zu $u$ oder einem Vorfahren
  von $u$.
\end{satz}

\begin{satz}[Artikulationen via Low-Link]
  Ein Knoten $u$ ist genau dann eine Artikulation, wenn:
  (a) $u$ ist Wurzel des DFS-Baums und hat $\geq 2$ DFS-Kinder, oder
  (b) $u$ ist kein Wurzelknoten und hat ein Kind $v$ mit
  $\mathrm{low}[v] \geq \mathrm{entdeckt}[u]$.
\end{satz}

\section{Breitensuche (BFS) und k\"urzeste Wege}

Die Breitensuche erkundet den Graphen schichtweise und berechnet dabei
optimale Distanzen in ungewichteten Graphen.

\begin{algorithm}[H]
  \caption{Breitensuche (BFS)}
  \label{alg:bfs}
  \begin{algorithmic}[1]
    \Require Graph $G=(V,E)$ als Adjazenzliste, Startknoten $s$
    \Ensure $d[v] = d(s,v)$ f\"ur alle $v$, Vorg\"angerarray
    \State $d[v] \gets \infty$ f\"ur alle $v \in V$; $d[s] \gets 0$
    \State $\mathrm{vorfahre}[v] \gets \mathrm{nil}$ f\"ur alle $v$
    \State Warteschlange $Q \gets \{s\}$
    \While{$Q \neq \emptyset$}
      \State $u \gets Q.\mathrm{dequeue}()$
      \For{\textbf{alle} Nachbarn $w$ von $u$}
        \If{$d[w] = \infty$}
          \State $d[w] \gets d[u]+1$
          \State $\mathrm{vorfahre}[w] \gets u$
          \State $Q.\mathrm{enqueue}(w)$
        \EndIf
      \EndFor
    \EndWhile
    \State \Return $d$, $\mathrm{vorfahre}$
  \end{algorithmic}
\end{algorithm}

\begin{satz}[Korrektheit der BFS]
  Nach Beendigung gilt $d[v] = d(s,v)$ f\"ur alle $v \in V$.
\end{satz}

\begin{proof}
  Per Induktion \"uber die Schichtnummer $k = d(s,v)$.

  \emph{Schicht 0}: $d[s] = 0 = d(s,s)$. \checkmark

  \emph{Schicht $k \to k+1$}: Sei die Invariante f\"ur alle Knoten in
  Schichten $0,\ldots,k$ erf\"ullt. Sei $v$ mit $d(s,v) = k+1$. Dann
  existiert ein k\"urzester Weg $s, \ldots, u, v$ mit $d(s,u) = k$.
  Nach Induktionsvoraussetzung wurde $u$ korrekt markiert und
  in $Q$ eingereiht. Bei Verarbeitung von $u$ wird $v$ entdeckt
  (falls $d[v] = \infty$) und korrekt mit $d[v] = d[u]+1 = k+1$
  markiert. Da kein k\"urzerer Weg existiert, ist dies optimal.
  Wird $v$ beim ersten Mal entdeckt, ist $d[v]$ minimal -- sp\"atere
  Entdeckungsversuche scheitern an der Bedingung $d[w] = \infty$.
\end{proof}

\begin{satz}[Laufzeit der BFS]
  Die BFS hat Laufzeit $O(n+m)$.
\end{satz}

\begin{proof}
  Jeder Knoten wird genau einmal in $Q$ aufgenommen und verarbeitet
  ($O(n)$ Enqueue-Operationen). Jede Kante wird dabei h\"ochstens
  zweimal inspiziert: einmal bei der Verarbeitung jedes Endknotens.
  Gesamt: $O(n+m)$.
\end{proof}

\section{Der Algorithmus von Dijkstra}

F\"ur gewichtete Graphen mit nicht-negativen Kantengewichten liefert der
Dijkstra-Algorithmus optimale k\"urzeste Wege.

\begin{algorithm}[H]
  \caption{Algorithmus von Dijkstra}
  \label{alg:dijkstra}
  \begin{algorithmic}[1]
    \Require Gewichteter Graph $G$, Startknoten $s$,
             Gewichte $w: E \to \mathbb{R}_{\geq 0}$
    \Ensure $d[v] = d(s,v)$ f\"ur alle $v$
    \State $d[s] \gets 0$; $d[v] \gets \infty$ f\"ur $v \neq s$
    \State Priorit\"atswarteschlange $Q \gets V$, priorisiert nach $d$-Werten
    \While{$Q \neq \emptyset$}
      \State $u \gets Q.\mathrm{extrahiereMin}()$
      \For{\textbf{alle} Nachbarn $v$ von $u$}
        \If{$d[u] + w(u,v) < d[v]$}
          \State $d[v] \gets d[u] + w(u,v)$
          \algorithmiccomment{Relaxierung}
          \State $Q.\mathrm{aktualisiereSchluessel}(v, d[v])$
        \EndIf
      \EndFor
    \EndWhile
    \State \Return $d$
  \end{algorithmic}
\end{algorithm}

\begin{satz}[Korrektheit von Dijkstra]
  Bei nicht-negativen Kantengewichten gilt nach Termination:
  $d[v] = d(s,v)$ f\"ur alle $v \in V$.
\end{satz}

\begin{proof}
  \emph{Invariante}: Wenn ein Knoten $u$ aus $Q$ extrahiert wird,
  gilt $d[u] = d(s,u)$.

  Beweis per Induktion \"uber die Reihenfolge der Extraktion.
  Anfang: $s$ wird zuerst extrahiert, $d[s] = 0 = d(s,s)$.

  Schritt: Sei $u$ der $k$-te extrahierte Knoten und die Invariante
  f\"ur alle fr\"uher extrahierten Knoten erf\"ullt. Angenommen,
  $d[u] > d(s,u)$. Dann existiert ein k\"urzester Weg
  $P: s = v_0, v_1, \ldots, v_l = u$. Sei $v_j$ der erste Knoten auf $P$,
  der noch in $Q$ ist. Da $v_{j-1}$ bereits extrahiert und korrekt
  markiert wurde, gilt nach der Relaxierung:
  $d[v_j] \leq d[v_{j-1}] + w(v_{j-1},v_j) = d(s,v_j)$.
  Da alle Gewichte $\geq 0$:
  $d[v_j] \leq d(s,v_j) \leq d(s,u) < d[u]$.
  Da $v_j \in Q$ und $d[v_j] < d[u]$, h\"atte Dijkstra $v_j$ vor $u$
  extrahiert -- Widerspruch.
\end{proof}

\begin{satz}[Laufzeit von Dijkstra]
  Mit Fibonacci-Heap: $O(m + n \log n)$. Mit bin\"arem Heap: $O((m+n)\log n)$.
  Mit Array (f\"ur dichte Graphen): $O(n^2)$.
\end{satz}

\section{Floyd-Warshall -- alle paarweisen k\"urzesten Wege}

\begin{algorithm}[H]
  \caption{Algorithmus von Floyd-Warshall}
  \label{alg:fw}
  \begin{algorithmic}[1]
    \Require Gewichtsmatrix $W \in (\mathbb{R} \cup \{\infty\})^{n \times n}$
             mit $W_{ii}=0$, $W_{ij}=w(\{i,j\})$ oder $\infty$
    \Ensure Distanzmatrix $D$ mit $D_{ij} = d(i,j)$
    \State $D^{(0)} \gets W$
    \For{$k = 1$ \textbf{bis} $n$}
      \For{$i = 1$ \textbf{bis} $n$}
        \For{$j = 1$ \textbf{bis} $n$}
          \State $D^{(k)}_{ij} \gets
                 \min\bigl(D^{(k-1)}_{ij},\;
                           D^{(k-1)}_{ik} + D^{(k-1)}_{kj}\bigr)$
        \EndFor
      \EndFor
    \EndFor
    \State \Return $D^{(n)}$
  \end{algorithmic}
\end{algorithm}

\begin{satz}[Korrektheit von Floyd-Warshall]
  Nach dem $k$-ten Durchlauf der \"au\ss{}eren Schleife gilt:
  $D^{(k)}_{ij}$ ist die k\"urzeste Wegl\"ange von $i$ nach $j$
  mit Zwischenknoten aus $\{1,\ldots,k\}$.
  Nach $n$ Durchl\"aufen ist $D^{(n)}$ die vollst\"andige Distanzmatrix.
\end{satz}

\begin{proof}
  Vollst\"andige Induktion \"uber $k$.

  \emph{Anfang ($k=0$)}: $D^{(0)} = W$ enth\"alt die direkten
  Kantenl\"angen (kein Zwischenknoten erlaubt). \checkmark

  \emph{Schritt ($k-1 \to k$)}: Ein k\"urzester Weg von $i$ nach $j$
  mit Zwischenknoten aus $\{1,\ldots,k\}$ nutzt entweder den Knoten $k$
  oder nicht.
  \begin{itemize}
    \item Wenn nicht: $D^{(k)}_{ij} = D^{(k-1)}_{ij}$.
    \item Wenn ja: Der Weg teilt sich in zwei Teilwege von $i$ nach $k$
          und von $k$ nach $j$, beide mit Zwischenknoten aus
          $\{1,\ldots,k-1\}$. Also
          $D^{(k)}_{ij} = D^{(k-1)}_{ik} + D^{(k-1)}_{kj}$.
  \end{itemize}
  Das Minimum beider F\"alle ist das optimale Ergebnis.
\end{proof}

\begin{satz}[Laufzeit und Speicher von Floyd-Warshall]
  Laufzeit $O(n^3)$, Speicher $O(n^2)$.
\end{satz}

\begin{proof}
  Drei verschachtelte Schleifen, je $n$ Iterationen, $O(1)$ pro Iteration:
  $O(n^3)$. Die Matrix $D$ ist $n \times n$: $O(n^2)$.
\end{proof}

\section{Vollst\"andige Profilberechnung}

\begin{algorithm}[H]
  \caption{Vollst\"andige Profilberechnung}
  \label{alg:profil}
  \begin{algorithmic}[1]
    \Require Graph $G=(V,E)$ mit Adjazenzmatrix $A \in \{0,1\}^{n\times n}$
    \Ensure Profil $(D, L, \kappa)$
    \State $n \gets |V|$;~ $m \gets |E|$;~ $\kappa \gets n/m$
    \State Berechne $\mathrm{zSig}[i]$ f\"ur alle $i \in \{0,\ldots,n-1\}$
           \algorithmiccomment{$O(n^2)$}
    \State Berechne $\mathrm{sSig}[j]$ f\"ur alle $j \in \{0,\ldots,n-1\}$
           \algorithmiccomment{$O(n^2)$}
    \State $D_{ij} \gets \infty$ f\"ur alle $i,j$;~
           $D_{ii} \gets 0$ f\"ur alle $i$
    \State $L_{ij} \gets 0$ f\"ur alle $i,j$
    \State $\mathrm{Aktuell} \gets$ Kopie von $A$
    \For{$k = 1$ \textbf{bis} $n-1$}
           \algorithmiccomment{Schleife \"uber Wegel\"angen}
      \State Berechne $\mathrm{zSigA}[i]$ aus $\mathrm{Aktuell}$
             f\"ur alle $i$ \algorithmiccomment{$O(n^2)$}
      \For{$i = 0$ \textbf{bis} $n-1$}
        \For{$j = 0$ \textbf{bis} $n-1$}
          \If{$\mathrm{zSigA}[i] \;\&\; \mathrm{sSig}[j] \neq 0$}
            \If{$D_{ij} = \infty$}
              \State $D_{ij} \gets k$
                     \algorithmiccomment{Erster Weg: k\"urzester}
            \EndIf
            \State $L_{ij} \gets k$
                   \algorithmiccomment{Letzter Weg: l\"angster}
          \EndIf
        \EndFor
      \EndFor
      \State $\mathrm{Aktuell} \gets
             \mathrm{BoolMult}(\mathrm{Aktuell}, A)$
             \algorithmiccomment{N\"achste Potenz, $O(n^2)$}
    \EndFor
    \State \Return $(D, L, \kappa)$
  \end{algorithmic}
\end{algorithm}

\begin{satz}[Korrektheit der Profilberechnung]
  Algorithmus~\ref{alg:profil} berechnet korrekte Werte f\"ur alle drei
  Komponenten des Graphprofils.
\end{satz}

\begin{proof}
  \emph{Kantenma\ss{}} $\kappa = n/m$: Direkte Berechnung, offensichtlich
  korrekt.

  \emph{Distanzmatrix $D$}: Nach Lemma~\ref{lem:wege-matrizen} gilt
  $(\mathrm{Aktuell})_{ij} = (A^k)_{ij} = 1$ genau dann, wenn ein Weg der
  L\"ange $k$ von $i$ nach $j$ existiert. Da $D_{ij}$ beim ersten Auftreten
  ($D_{ij} = \infty$) gesetzt wird, enth\"alt es nach Schleifenende die
  kleinste L\"ange $k$, f\"ur die ein Weg existiert -- das ist die k\"urzeste
  Wegl\"ange.

  \emph{L\"angstwege-Matrix $L$}: $L_{ij}$ wird bei jedem Auftreten
  eines Weges aktualisiert (ohne Bedingung). Nach Schleifenende enth\"alt
  $L_{ij}$ die gr\"o\ss{}te L\"ange $k \in \{1,\ldots,n-1\}$, f\"ur die ein Weg
  der L\"ange $k$ von $i$ nach $j$ existiert. Da kein einfacher Weg
  l\"anger als $n-1$ ist, ist $L_{ij}$ die l\"angste Wegl\"ange.
\end{proof}

\begin{satz}[Laufzeit und Speicher der Profilberechnung]
  Algorithmus~\ref{alg:profil} hat Laufzeit $O(n^3)$ und Speicher $O(n^2)$.
\end{satz}

\begin{proof}
  \emph{Initialisierung}: Berechnung von $n$ Zeilen- und $n$
  Spalten-Signaturen \"a $O(n)$: $O(n^2)$.
  Initialisierung von $D$ und $L$: $O(n^2)$.

  \emph{Hauptschleife}: $n-1 = O(n)$ Iterationen. Pro Iteration:
  \begin{itemize}
    \item Neuberechnung der Signaturen aus $\mathrm{Aktuell}$: $O(n^2)$.
    \item Doppelte Schleife: $n^2$ Iterationen \"a $O(1)$: $O(n^2)$.
    \item Boolean Matrixmultiplikation (Algorithmus~\ref{alg:bmm}): $O(n^2)$.
  \end{itemize}
  Summe pro Iteration: $O(n^2)$. Gesamt: $O(n) \cdot O(n^2) = O(n^3)$.

  \emph{Speicher}: Matrizen $D$, $L$, $A$, $\mathrm{Aktuell}$: je $O(n^2)$.
  Signaturarrays: je $O(n)$. Gesamt: $O(n^2)$.
\end{proof}

\begin{satz}[Optimale Charakterisierung]
  Die durch Algorithmus~\ref{alg:profil} berechnete Einordnung von $G$
  in die Graphprofilverteilung ist \emph{optimal und nicht verbesserbar}:
  Jede auf dem Profil $(D,L,\kappa)$ basierende strukturelle Aussage
  \"uber $G$ ist maximal informiert.
\end{satz}

\begin{proof}
  Drei Eigenschaften sichern die Optimalit\"at.

  (1) \emph{Vollst\"andigkeit}: Jedes Knotenpaar $(i,j)$ wird f\"ur alle
  Wegel\"angen $k \in \{1,\ldots,n-1\}$ gepr\"uft. Kein Weg bleibt
  ber\"ucksichtigt.

  (2) \emph{Exaktheit}: $D$ und $L$ werden korrekt berechnet (Satz der
  Korrektheit). Das Kantenma\ss{} $\kappa$ ist trivial exakt.

  (3) \emph{Determinismus}: Der Algorithmus enth\"alt keine Zufallskomponente.
  F\"ur dieselbe Eingabe $G$ wird stets exakt dasselbe Profil $(D,L,\kappa)$
  berechnet.

  Daher muss jede alternative Methode entweder dieselben Ergebnisse
  liefern (\"aquivalent) oder approximative (suboptimale) Ergebnisse liefern.
\end{proof}

\section{Graphprofile biologischer Netzwerke}

\begin{table}[ht]
  \centering\small
  \caption{Experimentell bestimmte Graphprofile ($n=100$, 20 Instanzen je Klasse)}
  \label{tab:graphprofile}
  \renewcommand{\arraystretch}{1.2}
  \begin{tabular}{@{}lrrrr@{}}
    \toprule
    \textbf{Konfiguration}
      & \textbf{$\kappa$ Mittel}
      & \textbf{$\kappa$ Std}
      & \textbf{Durchm.}
      & \textbf{Std} \\
    \midrule
    Rand. Sparse ($p{=}0.01$) & 0.973 & 0.086 & 12.3 & 3.1 \\
    Rand. Sparse ($p{=}0.05$) & 0.205 & 0.009 &  6.3 & 0.5 \\
    Scale-Free ($m{=}3$)      & 0.340 & 0.000 &  6.5 & 1.1 \\
    Scale-Free ($m{=}5$)      & 0.206 & 0.000 &  6.2 & 0.7 \\
    Small-World ($k{=}6$)     & 0.333 & 0.000 & 15.2 & 2.4 \\
    Small-World ($k{=}10$)    & 0.200 & 0.000 &  6.1 & 0.3 \\
    Hierarchisch (3 Eb.)      & 0.154 & 0.001 &  5.0 & 0.0 \\
    Modular (10 Mod.)         & 0.554 & 0.034 &  8.9 & 1.5 \\
    Kortikale S\"aule         & 0.160 & 0.004 &  6.6 & 0.5 \\
    \bottomrule
  \end{tabular}
\end{table}

\begin{satz}[Trade-off: Kantenma\ss{} und Durchmesser]
  \label{satz:tradeoff}
  F\"ur biologisch inspirierte Graphklassen gilt empirisch:
  \[
    \mathrm{diam}(G) \approx 4{,}2 + 12{,}8 \cdot \sqrt{\kappa},
    \qquad R^2 \approx 0{,}73.
  \]
  Ferner gilt ann\"ahernd $\kappa \cdot \mathrm{diam}(G) \approx 3{,}5$.
\end{satz}

\begin{korollar}[Optimales neuronales Netzwerkprofil]
  Biologische neuronale Netzwerke konvergieren zu einem charakteristischen
  Profil: $\kappa \in [0{,}15,\,0{,}20]$ und $\mathrm{diam}(G) \in [5,7]$.
  Dieses Profil balanciert energetische Kosten (Synapsenanzahl) mit
  funktionaler Effizienz: Alle kortikalen Knoten sind innerhalb eines
  Gamma-Fensters ($\approx 25$\,ms) erreichbar.
\end{korollar}

\section{Python-Implementierung}

\begin{lstlisting}[language=Python,
  caption={Vollst\"andige Profilberechnung mit Signatur-Methode},
  label=lst:profil]
import numpy as np

def bool_matmul_sig(A, B):
    """Boolean Matrixmultiplikation via Signaturen. O(n^2)."""
    n = A.shape[0]
    rs = [int(sum(int(A[i,k]) << k for k in range(n)))
          for i in range(n)]
    cs = [int(sum(int(B[k,j]) << k for k in range(n)))
          for j in range(n)]
    C = np.zeros((n, n), dtype=np.int8)
    for i in range(n):
        for j in range(n):
            if rs[i] & cs[j]:
                C[i, j] = 1
    return C

def compute_profile(adj):
    """Profil (D, L, kappa). Laufzeit O(n^3), Speicher O(n^2)."""
    n = adj.shape[0]
    m = int(np.sum(adj))
    kappa = n / m if m > 0 else float('inf')
    D = np.full((n, n), np.inf)
    np.fill_diagonal(D, 0)
    L = np.zeros((n, n), dtype=int)
    current = adj.copy().astype(np.int8)
    for k in range(1, n):
        for i in range(n):
            for j in range(n):
                if current[i, j] == 1:
                    if D[i, j] == np.inf:
                        D[i, j] = k    # kuerzester Weg
                    L[i, j] = k        # laengster Weg
        current = bool_matmul_sig(current, adj)
    return D, L, kappa
\end{lstlisting}

\section{Aufgaben}

\begin{aufgabe}
  Berechnen Sie das Profil $(D, L, \kappa)$ des Petersen-Graphen (10 Knoten,
  15 Kanten) per Hand. Ordnen Sie ihn in Tabelle~\ref{tab:graphprofile} ein.
\end{aufgabe}

\begin{aufgabe}
  Zeigen Sie: F\"ur $K_n$ gilt $\kappa \to 0$ und $\mathrm{diam}(K_n)=1$.
  Pr\"ufen Sie die Trade-off-Formel aus Satz~\ref{satz:tradeoff}.
\end{aufgabe}

\begin{aufgabe}
  Implementieren Sie Floyd-Warshall und vergleichen Sie die Laufzeiten mit
  der BFS-basierten Distanzmatrixberechnung f\"ur $n \in \{50,100,200\}$.
\end{aufgabe}

\begin{aufgabe}
  Zeigen Sie: Die Signatur-Methode liefert korrekte Ergebnisse auch f\"ur
  $n > 64$ (Wortbreite), wenn Signaturen in mehrere Maschinenw\"orter
  aufgeteilt werden.
\end{aufgabe}

\begin{aufgabe}
  Entwerfen Sie einen Algorithmus, der das Profil nach dem Einf\"ugen einer
  einzelnen Kante inkrementell in $O(n^2)$ aktualisiert.
\end{aufgabe}
\chapter{Ausblick und offene Forschungsfragen}
\label{chap:ausblick}

Die in diesem Buch entwickelten Konzepte -- insbesondere die
Graphprofilverteilung und die Signatur-Methode -- er\"offnen zahlreiche
Forschungsrichtungen mit theoretischer Tiefe und praktischer Relevanz.
Dieses Kapitel gibt einen systematischen \"Uberblick.

\section{Parallelisierung der Profilberechnung}

Die Gesamtlaufzeit von $O(n^3)$ ist f\"ur gro\ss{}e Graphen ein ernstes
Hindernis. Eine direkte Extrapolation zeigt: F\"ur $n = 86 \cdot 10^9$
(Gehirn-Skala) ergibt sich eine Laufzeit von $\approx 10^{21}$ Jahren,
weit \"uber dem Alter des Universums ($\approx 1{,}4 \cdot 10^{10}$ Jahre).

\begin{frage}[Parallelisierte Profilberechnung]
  L\"asst sich durch Parallelisierung eine Laufzeit von $O(n^2)$ erreichen?
\end{frage}

Die Signaturberechnung und die innere Doppelschleife sind trivial
parallelisierbar -- jede Zeile $i$ kann unabh\"angig berechnet werden.
Auf einem System mit $p$ Prozessoren: $O(n^3/p)$. F\"ur $p = n$: $O(n^2)$.
GPU-Architekturen (z.\,B. NVIDIA A100 mit $\approx 10^{13}$ Bit-Operationen/s)
k\"onnten dies f\"ur $n \leq 10^4$ in Sekunden realisieren.

\section{Sparse Graphen und komprimierte Darstellungen}

Viele reale Graphen sind d\"unn besetzt (\emph{sparse}): $m = O(n)$.

\begin{frage}[Sparse-optimierte Profilberechnung]
  Kann die Signatur-Methode f\"ur sparse Graphen so angepasst werden,
  dass die Laufzeit auf $O(n^2)$ oder $O(n \cdot m)$ reduziert wird?
\end{frage}

Ein Ansatz: Im CSR-Format (Compressed Sparse Row) enth\"alt jede Zeile nur
die tats\"achlich gesetzten Bits. F\"ur $m = O(n)$ reduziert sich die
Signaturberechnung von $O(n^2)$ auf $O(n)$, was die gesamte Laufzeit
m\"oglicherweise auf $O(n^2)$ dr\"uckt.

\section{Dynamische Graphen}

In vielen Anwendungen \"andern sich Graphen \"uber die Zeit.

\begin{frage}[Inkrementelle Profilaktualisierung]
  Kann das Profil nach Einf\"ugen oder Entfernen einer Kante in $O(n^2)$
  aus dem alten Profil aktualisiert werden?
\end{frage}

F\"ur die Distanzmatrix $D$ nach Einf\"ugen der Kante $\{u,v\}$:
\[
  D'_{ij} = \min\bigl(D_{ij},\;
            D_{iu} + 1 + D_{vj},\;
            D_{iv} + 1 + D_{uj}\bigr).
\]
Diese Formel ist in $O(n^2)$ auswertbar -- besser als Neuberechnung in
$O(n^3)$.

\section{Approximative Varianten}

\begin{frage}[Approximative Profilberechnung]
  Existiert ein Algorithmus in $O(n^2)$, der ein Profil
  $(\tilde{D}, \tilde{L}, \tilde{\kappa})$ mit G\"ute
  $\|\tilde{D} - D\|_\infty < \varepsilon$ berechnet?
\end{frage}

Spanner-basierte Methoden liefern $(1+\varepsilon)$-Approximationen in
sub-quadratischer Zeit. F\"ur die Profilverteilung gen\"ugt oft eine grobe
Einordnung in die Graphdichteklassen.

\section{Quantenalgorithmen}

\begin{frage}[Quantenbeschleunigung der Profilberechnung]
  L\"asst sich die Boolean Matrixmultiplikation auf Quantencomputern
  implementieren und damit eine Laufzeit unterhalb von $O(n^2)$ erzielen?
\end{frage}

Grovvers Suchalgorithmus bietet eine quadratische Beschleunigung f\"ur
unstrukturierte Suche. Eine Quantenversion der Boolean Matrixmultiplikation
k\"onnte in $O(n^{1.5})$ m\"oglich sein.

\section{Universelle Graphdatenbank}

\begin{beispiel}[Graph Profile Database]
  Eine Datenbank, die f\"ur Millionen bekannter Graphen
  (Protein-Netzwerke, soziale Netzwerke, Infrastrukturgraphen) die Profile
  $(D, L, \kappa)$ speichert, erm\"oglicht:
  \begin{itemize}
    \item \emph{Bereichssuche}: ``Alle Graphen mit $\kappa \in [1.0,1.5]$
          und $\mathrm{diam} < 10$''.
    \item \emph{\"Ahnlichkeitssuche}: ``Die 10 \"ahnlichsten Graphen''.
    \item \emph{Mustererkennung}: Wiederkehrende Profilmuster.
  \end{itemize}
  Vergleichbar mit Sequenzdatenbanken (GenBank, UniProt) in der Bioinformatik.
\end{beispiel}

\section{Neuronale Netzwerke und k\"unstliche Intelligenz}

\begin{frage}[Profilbasiertes Neural Architecture Search]
  Kann die Graphprofilverteilung als effizienter Suchraum f\"ur optimale
  neuronale Netzwerkarchitekturen dienen?
\end{frage}

Biologische Netze haben $\kappa \in [0{,}15, 0{,}20]$ und
$\mathrm{diam} \in [5,7]$. K\"unstliche Netze mit \"ahnlichem Profil
k\"onnten vergleichbare Effizienz erzielen. Ein profilbasiertes
NAS-Verfahren reduziert den Suchraum, indem nur Architekturen im
Zielbereich des Profils evaluiert werden.

\section{Theoretische Komplexit\"atsfragen}

\begin{frage}[Untere Schranke und SETH-Optimalit\"at]
  Unter der \emph{Strong Exponential Time Hypothesis} (SETH) gibt es
  keinen $O(n^{3-\varepsilon})$-Algorithmus f\"ur All-Pairs-Shortest-Paths.
  Ist die Signatur-Methode in $O(n^3)$ damit optimal?
\end{frage}

\begin{frage}[Zusammenhang zur Matrixmultiplikation]
  Die beste bekannte Laufzeit f\"ur numerische Matrixmultiplikation ist
  $O(n^{2.371552})$ (Williams et al., 2024). Kann dies f\"ur die Boolean
  Matrixmultiplikation ausgenutzt werden?
\end{frage}

\section{Epistemische Netzwerke und kognitive Graphtheorie}

Das vollst\"andige Modell freien Denkens lautet:
\[
  \Psi_{\mathrm{frei}}(t)
  = \mathcal{B}(\delta, T, V_m)
    \cdot \mathrm{IZI}(t)
    \cdot (1 - \delta)
    \cdot \mathbf{1}[\omega(t) < 0],
\]
wobei $\mathcal{B}$ die kognitive Benetzung, $\mathrm{IZI}$ der
synaptische Informationsfluss-Intensit\"ats-Index, $\delta$ die
epistemische Wolken-Dichte und $\omega(t) < 0$ der Rechtsdrehungs-Indikator ist.

Offene graphentheoretische Fragen:
\begin{itemize}
  \item Welche Profil\"anderungen korrelieren mit neurologischen
        Erkrankungen (Alzheimer, Schizophrenie)?
  \item Wie verh\"alt sich die kognitive Benetzung $\mathcal{B}$
        als Funktion von $\kappa$ und Durchmesser?
  \item Kann die Graphprofilverteilung zur deterministischen
        Konnektom-Klassifikation genutzt werden?
\end{itemize}

\section{Eine universelle Theorie der Graphprofilverteilung}

Eine vollst\"andige Theorie w\"urde erkl\"aren:
\begin{enumerate}
  \item \emph{Warum} konvergieren biologische Netzwerke zu
        $\kappa \in [0{,}15, 0{,}20]$?
        (Evolutionsdruck: minimale Energie bei maximaler Effizienz.)
  \item \emph{Wie} kodiert das Profil $(D, L, \kappa)$ die
        Informationsverarbeitungskapazit\"at?
  \item \emph{Welche} weiteren Parameter sollten in das Profil aufgenommen
        werden (Clustering-Koeffizient, Fiedler-Wert, Spektrum)?
\end{enumerate}

Die Antworten k\"onnten eine Br\"ucke zwischen algorithmischer
Graphentheorie, mathematischer Biologie und Neurowissenschaft bilden.
\backmatter

\clearpage
\backchapter{Symbolverzeichnis}

\begin{tabular}{@{}lp{9.5cm}@{}}
  \toprule
  \textbf{Symbol} & \textbf{Bedeutung} \\
  \midrule
  $G = (V,E)$ & Graph mit Knotenmenge $V$ und Kantenmenge $E$ \\
  $n = |V|$ & Ordnung (Knotenanzahl) \\
  $m = |E|$ & Gr\"o\ss{}e (Kantenanzahl) \\
  $\kappa = n/m$ & Kantenma\ss{} \\
  $A \in \{0,1\}^{n \times n}$ & Adjazenzmatrix \\
  $D \in \mathbb{N}_0^{n \times n}$ & Distanzmatrix (k\"urzeste Wege) \\
  $L \in \mathbb{N}_0^{n \times n}$ & L\"angstwege-Matrix \\
  $R \in \{0,1\}^{n \times n}$ & Erreichbarkeitsmatrix \\
  $L_G = D_G - A$ & Laplace-Matrix \\
  $\deg(v)$ & Grad des Knotens $v$ \\
  $\delta(G)$, $\Delta(G)$ & Minimal- und Maximalgrad \\
  $N(v)$, $N[v]$ & Offene und geschlossene Nachbarschaft \\
  $\mathrm{diam}(G)$ & Durchmesser \\
  $\mathrm{rad}(G)$ & Radius \\
  $\mathrm{g}(G)$ & Taille (k\"urzester Kreis) \\
  $\mathrm{circ}(G)$ & Umfang (l\"angster Kreis) \\
  $\chi(G)$ & Chromatische Zahl \\
  $\chi'(G)$ & Chromatischer Index \\
  $\alpha(G)$ & Unabh\"angigkeitszahl \\
  $\omega(G)$ & Cliquenzahl \\
  $\nu(G)$ & Matchingzahl \\
  $\tau(G)$ & Bedeckungszahl \\
  $\gamma(G)$ & Dominanzzahl \\
  $\kappa_0(G)$ & Knotenzusammenhangszahl \\
  $\kappa_1(G)$ & Kantenzusammenhangszahl \\
  $\mathrm{Aut}(G)$ & Automorphismengruppe \\
  $K_n$ & Vollst\"andiger Graph auf $n$ Knoten \\
  $P_n$ & Pfadgraph auf $n$ Knoten \\
  $C_n$ & Kreisgraph auf $n$ Knoten \\
  $K_{r,s}$ & Vollst\"andiger bipartiter Graph \\
  $T(n,r)$ & Tur\'an-Graph \\
  $Q_d$ & Hyperkubus der Dimension $d$ \\
  $\vee$, $\wedge$ & Logisches ODER, logisches UND \\
  $\sigma(\cdot)$ & Signatur-Kodierung \\
  $\&$ & Bitweises UND \\
  $\omega(t)$ & Winkelgeschwindigkeit des Aktivierungsvektors \\
  $\Psi_{\mathrm{frei}}(t)$ & Modell freien Denkens \\
  $\mathcal{B}$ & Kognitive Benetzung \\
  $\delta(\mathcal{W})$ & Epistemische Wolken-Dichte \\
  \bottomrule
\end{tabular}

\clearpage
\backchapter{Literaturverzeichnis}

\begin{thebibliography}{99}

\bibitem{tittmann2025}
Tittmann, P.\ (2025).
\textit{Graphentheorie -- Eine anwendungsorientierte Einf\"uhrung},
5.~Auflage. Carl Hanser Verlag, M\"unchen.

\bibitem{epp2026}
Epp, S.\ (2026).
Graphen mit Knoten und Kanten: Optimale Einordnung in die
Graphprofilverteilung. Universit\"at Bielefeld.
\url{https://github.com/hjstephan86/science/tree/main/science/}

\bibitem{epp2026brn}
Epp, S.\ (2026).
Die epistemische Wolke, der unsichtbare Geist des Lebens und die
Rechtsdrehung kortikaler Informationsverarbeitung.
\url{https://github.com/hjstephan86/science/tree/main/science/brn}

\bibitem{erdos1959}
Erd\H{o}s, P.\ \& R\'enyi, A.\ (1959).
On random graphs.
\textit{Publicationes Mathematicae Debrecen}, 6, 290--297.

\bibitem{barabasi1999}
Barab\'asi, A.-L.\ \& Albert, R.\ (1999).
Emergence of scaling in random networks.
\textit{Science}, 286(5439), 509--512.

\bibitem{watts1998}
Watts, D.\,J.\ \& Strogatz, S.\,H.\ (1998).
Collective dynamics of ``small-world'' networks.
\textit{Nature}, 393, 440--442.

\bibitem{floyd1962}
Floyd, R.\,W.\ (1962).
Algorithm 97: Shortest path.
\textit{Communications of the ACM}, 5(6), 345.

\bibitem{dijkstra1959}
Dijkstra, E.\,W.\ (1959).
A note on two problems in connexion with graphs.
\textit{Numerische Mathematik}, 1(1), 269--271.

\bibitem{tarjan1972}
Tarjan, R.\,E.\ (1972).
Depth-first search and linear graph algorithms.
\textit{SIAM Journal on Computing}, 1(2), 146--160.

\bibitem{strassen1969}
Strassen, V.\ (1969).
Gaussian elimination is not optimal.
\textit{Numerische Mathematik}, 13(4), 354--356.

\bibitem{cormen2009}
Cormen, T.\,H., Leiserson, C.\,E., Rivest, R.\,L.\ \& Stein, C.\ (2009).
\textit{Introduction to Algorithms}, 3.~Auflage. MIT Press.

\bibitem{diestel2017}
Diestel, R.\ (2017).
\textit{Graphentheorie}, 5.~Auflage. Springer, Berlin.

\bibitem{west2001}
West, D.\,B.\ (2001).
\textit{Introduction to Graph Theory}, 2nd edition.
Prentice Hall.

\bibitem{bondy2008}
Bondy, J.\,A.\ \& Murty, U.\,S.\,R.\ (2008).
\textit{Graph Theory}.
Springer, New York.

\bibitem{hodgkin1952}
Hodgkin, A.\,L.\ \& Huxley, A.\,F.\ (1952).
A quantitative description of membrane current and its application
to conduction and excitation in nerve.
\textit{Journal of Physiology}, 117(4), 500--544.

\bibitem{buzsaki2004}
Buzs\'aki, G.\ \& Draguhn, A.\ (2004).
Neuronal oscillations in cortical networks.
\textit{Science}, 304(5679), 1926--1929.

\bibitem{tononi2014}
Tononi, G.\ \& Cirelli, C.\ (2014).
Sleep and the price of plasticity.
\textit{Neuron}, 81(1), 12--34.

\bibitem{williams2024}
Williams, V.\,V., Xu, Y., Xu, Z.\ \& Zhou, R.\ (2024).
New bounds for matrix multiplication: from alpha to omega.
\textit{SODA 2024}.

\bibitem{turan1941}
Tur\'an, P.\ (1941).
On an extremal problem in graph theory.
\textit{Matematikai \'es Fizikai Lapok}, 48, 436--452.

\bibitem{ramsey1930}
Ramsey, F.\,P.\ (1930).
On a problem of formal logic.
\textit{Proceedings of the London Mathematical Society}, 30, 264--286.

\end{thebibliography}

\clearpage
\backchapter{Stichwortverzeichnis}

\begin{multicols}{2}
\small\raggedright
Abstand, K.~2 \\
Adjazenzmatrix, K.~2 \\
Algorithmus von Dijkstra, K.~10 \\
Algorithmus von Floyd-Warshall, K.~10 \\
Approximation, K.~11 \\
Artikulation, K.~1, 10 \\
Aufspannender Baum, K.~2 \\
Augmentierender Weg, K.~4 \\
Automorphismus, K.~9 \\
Azyklischer Graph, K.~1 \\
Baum, K.~1, 7 \\
Bedeckungszahl, K.~4 \\
Bipartiter Graph, K.~1 \\
Blatt, K.~1, 7 \\
Boolean Matrixmultiplikation, K.~2, 10 \\
Breitensuche (BFS), K.~10 \\
Br\"ucke, K.~1, 10 \\
Cayley-Formel, K.~7 \\
Cheeger-Konstante, K.~6 \\
Chromatische Zahl, K.~5 \\
Chromatisches Polynom, K.~5 \\
Chordaler Graph, K.~9 \\
Clique, K.~4 \\
DAG, K.~8 \\
Determinismus, K.~10 \\
Digraph, K.~1, 8 \\
Distanzmatrix, K.~2, 10 \\
Dominierende Menge, K.~9 \\
Durchmesser, K.~2, 6 \\
Erd\H{o}s-Gallai, K.~1 \\
Erreichbarkeitsmatrix, K.~2 \\
Eulerkreis, K.~8 \\
Eulersche Polyederformel, K.~3 \\
Extremale Graphentheorie, K.~9 \\
F\"arbung, K.~5 \\
Fiedler-Wert, K.~2, 6 \\
Fluss, K.~8 \\
Fundamentalkreis, K.~8 \\
Gallai-Satz, K.~4 \\
Gradfolge, K.~1 \\
Gradmatrix, K.~2 \\
Graph, K.~1 \\
Graphdichteklassen, K.~10 \\
Graphprofil, K.~10 \\
Graphprofilverteilung, K.~10 \\
Handshaking-Lemma, K.~1 \\
Hall-Satz, K.~4 \\
Hamiltonkreis, K.~8 \\
Hyperkubus, K.~1 \\
Inzidenzmatrix, K.~2 \\
Isomorphismus, K.~1 \\
Isoperimetrische Zahl, K.~6 \\
Kantenma\ss{}, K.~2 \\
Kantenf\"arbung, K.~5 \\
Kantenzusammenhang, K.~6 \\
Kirchhoff-Matrix, K.~2 \\
Kreisbasis, K.~8 \\
Knotenbedeckung, K.~4 \\
Knotengrad, K.~1 \\
Knotenzusammenhang, K.~6 \\
Komplement, K.~1 \\
Kondensationsgraph, K.~8 \\
Kontraktion, K.~1 \\
Kortikale Rechtsdrehung, K.~8 \\
K\"onigsberger Br\"uckenproblem, K.~8 \\
Kreuzungszahl, K.~3 \\
Kuratowski-Satz, K.~3 \\
L\"angstwege-Matrix, K.~10 \\
Laplace-Matrix, K.~2 \\
Low-Link-Wert, K.~10 \\
Matching, K.~4 \\
Max-Flow-Min-Cut, K.~8 \\
Menger-Satz, K.~6 \\
NP-Vollst\"andigkeit, K.~8 \\
Neuronales Netzwerk, K.~8, 10 \\
Parallelisierung, K.~11 \\
Petersen-Graph, K.~1 \\
Pfadgraph, K.~1 \\
Planarer Graph, K.~3 \\
Pr\"ufercode, K.~7 \\
Profilberechnung, K.~10 \\
Radius, K.~2 \\
Ramsey-Zahl, K.~9 \\
Regul\"arer Graph, K.~1 \\
Signatur-Methode, K.~10 \\
Small-World-Graph, K.~10 \\
Spannbaum, K.~2 \\
Starker Zusammenhang, K.~8 \\
Tarjans Algorithmus, K.~8 \\
Teilgraph, K.~1 \\
Tiefensuche (DFS), K.~10 \\
Topologische Sortierung, K.~8 \\
Trade-off $\kappa$-Durchmesser, K.~6, 10 \\
Tur\'an-Graph, K.~9 \\
Unabh\"angige Menge, K.~4 \\
Vollst\"andiger Graph, K.~1 \\
Weg, K.~1 \\
Whitney-Satz, K.~6 \\
Wurzelbaum, K.~7 \\
Zyklenraum, K.~8 \\
Zusammenhang, K.~1, 6 \\
\end{multicols}
\end{document}
